\newfontfamily\BBSixGreek[Path=vendor/fonts/,BoldFont=DejaVuSerif-Bold.ttf,ItalicFont=DejaVuSerif.ttf,BoldItalicFont=DejaVuSerif-Bold.ttf]{DejaVuSerif.ttf}
\chapter{THÉORIES DU SIGNE : PEIRCE ET SAUSSURE}
\label{chapter-6}
\markright{Théories du signe : Peirce et Saussure}
\origpage{164}{150}
\begin{keybox}[Plan du chapitre]
\small
\textbf{I.\quad SAUSSURE ET PEIRCE : « DEUX GÉNIES ANTITHÉTIQUES »}
\begin{enumerate}[label=\arabic*.]
\item Différences entre sémiotique et sémiologie
\begin{enumerate}[label=\textcircled{\scriptsize\arabic*}]
\item Pragmatisme contre historicisme
\item Dyade contre triade
\end{enumerate}
\item Peirce : « le plus grand penseur américain qu’il y eut »
\item Saussure et l’acte de naissance de la sémiologie
\begin{enumerate}[label=\textcircled{\scriptsize\arabic*}]
\item Le \textit{Cours de linguistique générale} (1916)
\item L’acte de naissance de la sémiologie
\item La linguistique, branche et « patron » de la sémiologie
\end{enumerate}
\end{enumerate}
\textbf{II.\quad LA SÉMIOTIQUE TRIADIQUE DE PEIRCE}
\begin{enumerate}[label=\arabic*.]
\item Le modèle triadique
\item Les trois catégories phanéroscopique\ednote{原文如此；修饰复数 catégories 的形容词应为 phanéroscopiques。正文相应小节标题已用复数。}
\begin{enumerate}[label=\textcircled{\scriptsize\arabic*}]
\item La notion de \textit{phaneron}
\begin{enumerate}[label=\arabic*.]
\item La Priméité
\item La Secondéité
\item La Tiercéité
\end{enumerate}
\end{enumerate}
\item La tripartition du signe
\begin{enumerate}[label=\textcircled{\scriptsize\arabic*}]
\item Définitions du signe
\item Le \textit{representamen} (« signe »)
\item L’objet
\item L’« interprétant »
\item Illustration
\item La sémiosis : le processus de signification
\end{enumerate}
\item La trichotomie du signe
\begin{enumerate}[label=\textcircled{\scriptsize\arabic*}]
\item La trichotomie du representamen
\begin{enumerate}[label=\arabic*.]
\item Le qualisigne
\item La sinsigne\ednote{原文如此；应为 Le sinsigne，正文相应标题也使用 Le。}
\item Le légisigne
\end{enumerate}
\origpage{165}{151}
\item La trichotomie de l’objet
\begin{enumerate}[label=\arabic*.]
\item L’icone
\item L’indice
\item Le symbole
\item Le signal
\end{enumerate}
\item La trichotomie de l’interprétant
\begin{enumerate}[label=\arabic*.]
\item Le rhème
\item Le dicisigne
\item L’argument
\end{enumerate}
\end{enumerate}
\item Charles W. Morris : le successeur
\end{enumerate}
\textbf{III.\quad LE SIGNE DYADIQUE DE SAUSSURE}
\begin{enumerate}[label=\arabic*.]
\item Caractéristiques du signe linguistique
\begin{enumerate}[label=\textcircled{\scriptsize\arabic*}]
\item Dualité du signe linguistique : signifiant et signifié
\item Arbitraire du signe linguistique
\begin{enumerate}[label=\arabic*.]
\item L’arbitraire du signe chez Platon : le \textit{Cratyle}
\begin{enumerate}[label=\arabic*)]
\item La Thèse d’Hermogène
\item La Thèse de Cratyle
\item La Synthèse de Socrate
\end{enumerate}
\item L’arbitraire selon Saussure
\item Le cas des onomatopées
\end{enumerate}
\item La linéarité du signe linguistique
\item La notion de valeur linguistique
\begin{enumerate}[label=\arabic*.]
\item Chez Patañjali
\item Chez Saussure
\end{enumerate}
\item Immutabilité et mutabilité du signe linguistique
\begin{enumerate}[label=\arabic*.]
\item Immutabilité du signe
\item Mutabilité du signe
\end{enumerate}
\end{enumerate}
\item Les dichotomies du CLG
\begin{enumerate}[label=\textcircled{\scriptsize\arabic*}]
\item Linguistique interne et linguistique externe
\item Linguistique de la langue et linguistique de la parole
\begin{enumerate}[label=\arabic*.]
\item Le langage, hétéroclite et insaisissable
\item La langue, principe de classification
\item La parole, partie individuelle du langage
\item Langue ou parole, il faut choisir
\end{enumerate}
\item Linguistique synchronique et linguistique diachronique
\begin{enumerate}[label=\arabic*.]
\item La synchronie
\item La diachronie
\end{enumerate}
\origpage{166}{152}
\item Rapports syntagmatiques et rapports paradigmatiques
\begin{enumerate}[label=\arabic*.]
\item Les rapports syntagmatiques
\begin{enumerate}[label=\arabic*)]
\item Illustration \# 1
\item Illustration \# 2
\end{enumerate}
\item Les rapports paradigmatiques
\end{enumerate}
\end{enumerate}
\end{enumerate}
\textbf{Exercices}\par
\textbf{Pour aller plus loin}\par
\textbf{Glossaire}
\end{keybox}

\origpage{167}{153}
\begin{keybox}[Introduction]
Issu d’une famille genevoise d’illustres savants, Ferdinand de Saussure (1857-1913) a commencé sa carrière à Paris avec l’enseignement de grammaire comparée à l’École Pratique des Hautes Études de 1881 à 1891. Son célèbre \textit{Cours de linguistique générale} (1916), ouvrage fondateur de la linguistique moderne, est une œuvre posthume. Saussure a également défini les bases de la sémiologie, science des signes et de leur signification. Aux États-Unis, son contemporain Charles Sanders Peirce (1839-1914) logicien, philosophe et mathématicien de génie, travaillait lui aussi de son côté à l’édification de cette nouvelle science qu’il a nommée \textit{semiotics}.

Pour autant, il n’est pas tout à fait exact de dire que Saussure ou Peirce ont « inventé » la sémiologie. John Locke a été le premier à utiliser le terme de \textit{semeiotike}, à partir du grec ancien {\BBSixGreek σῆμα}\textit{sêma} qui signifie « signe ». Près d’un siècle plus tard, en 1764, c’est le philosophe et mathématicien Jean Henri Lambert (1728-1777) qui, s’inspirant de Locke, a développé dans la troisième partie de son \textit{Neues Organon} (« nouvel Organon ») une théorie générale des signes nommée \textit{semiotic}. Quoi qu’il en soit, c’est bien à Saussure et Peirce que revient le mérite d’avoir porté cette nouvelle science sur les fonds baptismaux.\ednote{原文 fonds baptismaux 如此；固定搭配应为 fonts baptismaux（洗礼盆），此处借指一门学科的创立。}
\end{keybox}

\section{SAUSSURE ET PEIRCE : « DEUX GÉNIES ANTITHÉTIQUES »}
En 1969, dans un article publié dans la revue \textit{Semiotica} créée la même année, \textbf{Émile Benveniste} présente le duo fondateur de la « science des signes », Peirce et Saussure, comme « deux génies antithétiques ». Pour lui, « ces deux esprits solitaires et singuliers » ont en commun de vouloir reprendre le projet du philosophe anglais \textbf{John Locke} (1632-1704) d’une théorie générale des signes. Benveniste ajoute que « leur formation, leur méthode, leur relation à l’objet de leur recherche diffèrent du tout au tout ». Avant de présenter les deux géants, quelques précisions terminologiques s’imposent.

\subsection{Différences entre sémiotique et sémiologie}
Si sémiologie et sémiotique peuvent être toutes deux définies comme l’étude des signes (linguistiques et non linguistiques), une question fréquemment posée est celle de la différence entre la sémiologie et la sémiotique. On répond en général qu’il n’y en pas et que ces deux termes renvient à la même science.\ednote{原文 « il n’y en pas »、« renvient » 如此；分别应为 « il n’y en a pas »、« renvoient »。} Ce n’est pas faux... mais ce n’est pas tout à fait vrai non plus.

Si sémiotique et sémiologie désignent, c’est vrai, la même discipline, c’est le terme « sémiotique » qui a été préféré en 1969 lors du premier congrès de l’Association Internationale de Sémiotique.\ednote{IASS–AIS 官方历届大会名录将 1974 年 Milano 大会列为第一届；因此不能称为 1969 年的首届大会。见\ecite{35}。} Donc, \textit{a priori}, la question est tranchée. Néanmoins, « sémiologie » continue à être utilisé, principalement en France. Voilà pour ce qui est de la terminologie. Au-delà de ces deux dénominations, il existe d’autres petites différences.

\origpage{168}{154}
\minititle{\textcircled{\scriptsize 1}\quad Pragmatisme contre historicisme}
D’abord, sémiotique et sémiologie ont chacune un « papa » différent, à savoir \textbf{Charles Sanders Peirce} aux États-Unis et \textbf{Ferdinand de Saussure} en Europe. La sémiotique est donc marquée par le « pragmatisme », courant philosophique typiquement américain qui postule qu’il faut « chercher le critère de la vérité dans l’action (\textit{pragma}) » et qu’« une idée n’est pas utile parce qu’elle est vraie, elle est vraie parce qu’elle est utile ». Cette philosophie fonde la vérité des idées non pas sur leur adéquation avec une quelconque réalité apparente ou cachée mais sur le fait qu’elles débouchent sur une action, ce qui finalement permet de juger de leur utilité.\ednote{此概括不宜直接当作 Peirce 的真理论。Peirce 的实用主义准则借可设想的实践效果澄清概念；真理则与无限探究所趋向的共同意见及独立于个体意见的实在相联，不等于“有用所以为真”。见 Peirce, « How to Make Our Ideas Clear » (1878)，\ecite{36}。} Saussure, lui, est européen et linguiste. En tant que penseur, il est imprégné par l’atmosphère intellectuelle de l’époque : d’une part, l’historicisme qui explique les faits humains par leur évolution, et d’autre part le positivisme qui proclame la rupture avec toute explication métaphysique en matière de connaissance. Ces contextes intellectuels et philosophiques expliquent donc la différence entre la sémiotique de Peirce et la sémiologie de Saussure.

\minititle{\textcircled{\scriptsize 2}\quad Dyade contre triade}
Autre différence (de taille, celle-ci) entre la sémiotique de Peirce et la sémiologie de Saussure : la sémiotique peircienne fonctionne sur une conception du signe \textit{triadique} (composé de trois éléments) alors que la sémiologie repose depuis Saussure sur un système \textit{dyadique} (composé de deux éléments). De nos jours, les concepts triadiques dominent dans les écoles sémiotiques anglo-saxonnes, tandis que les sémiologues européens (comme Roland Barthes, Louis Hjelmslev, Umberto Eco) sont généralement restés fidèles au dyadisme saussurien\origfoot{1}{On verra cela de façon approfondie dans les pages qui suivent ainsi que dans le chapitre 11, consacré à la sémantique.}.

\subsection{Peirce : « le plus grand penseur américain qu’il y eut »}
Dans un article de 1969 consacré à Saussure et Peirce, Émile Benveniste rappelle que Peirce est avant tout un savant : logicien, mathématicien, historien et philosophe des sciences. Ses travaux sur les fondements des mathématiques ont fait de Peirce un précurseur de \textbf{Bertrand Russell}. Ce dernier, dans \textit{Wisdom of the West}, (1959) lui a rendu hommage appuyé :
\begin{lecture}{}
Il ne fait aucun doute qu’il [Peirce] fut un des esprits les plus originaux du dix-neuvième siècle, et certainement le plus important penseur américain qu’il y eut.
\end{lecture}
Entré à 16 ans à l’Université d’Harvard dans laquelle son père était professeur de mathématiques, Peirce en est sorti diplômé en lettres (1862) et en chimie (1863). On raconte qu’à seize ans, le jeune Peirce a étudié \textit{La critique de la raison pure} de Kant et en a conclu que cet ouvrage était basé sur une « logique puérile ». Arrivé à l’âge adulte, en tant que savant et logicien, il s’est consacré à la recherche des conditions de vérité des assertions à caractère scientifique et, au-delà, à l’élaboration d’une théorie des conditions de la connaissance. En ce sens, sa \textit{semeiotic} est assimilée à la logique :
\origpage{169}{155}
\begin{lecture}{}
La logique, dans son sens général n’est qu’un autre nom de la sémiotique [semeiotic], la doctrine quasi-nécessaire ou formelle des signes.
\end{lecture}
Dans ce qui vient d’être dit, on voit que pour Peirce la sémiotique (tout comme la logique qui prête son langage formel à toutes les sciences) est une méta-science : en retrait par rapport aux autres sciences, elle les « coiffe » toutes en quelque sorte. Elle est ainsi, comme il le dit lui-même, un cadre de référence qui englobe toute étude.

Les principaux écrits de Peirce (près de 100 000 pages manuscrites) ont été vendus par sa femme en 1914 à l’Université de Harvard avant d’être recueillis dans des \textit{Collected Papers}. Depuis 1982, d’autres volumes ont commencé à être publiés dans le cadre d’une édition chronologique qui atteindra une trentaine de volumes, et dont 8 volumes sont déjà parus. On n’a donc pas fini de redécouvrir l’œuvre de ce génie prolifique.

\subsection{Saussure, et l’acte de naissance de la sémiologie}
De Saussure, on peut dire bien des choses, mais pas qu’il a été prolifique. En effet, à l’exception du mémoire et de la thèse mentionnés dans l’introduction, il n’a en effet à peu près rien n’écrit ni publié de son vivant.\ednote{原文 « rien n’écrit » 如此；应为 « rien écrit »。关于著述数量，\textit{Recueil des publications scientifiques de Ferdinand de Saussure} (1922) 除早期专著外还收录多篇论文，故“几乎没有再写作或发表”也过强；不能由 CLG 为身后编成推得这一结论。见\ecite{39}。}

\minititle{\textcircled{\scriptsize 1}\quad Le \textit{Cours de linguistique générale} (1916)}
Le \textit{Cours de Linguistique Générale}, œuvre posthume on le sait, a été rédigé par ses collègues avec l’aide d’un étudiant, en se fondant sur les notes de Saussure ainsi que celles prises par des étudiants lors des trois leçons de linguistique générale données en 1906-1907, 1908-1909 et 1910-1911.\ednote{此处 trois leçons 宜理解为三轮课程，而非三节课。CLG 第一版序言署 Charles Bally、Albert Sechehaye，书名页注明 Albert Riedlinger 协作；编辑者说明主要材料来自听课笔记，Saussure 留下的课程相关笔记非常有限。见\ecite{38}，第一版序言。}

\minititle{\textcircled{\scriptsize 2}\quad L’acte de naissance de la sémiologie}
Le \textit{Cours de linguistique générale} de Ferdinand de Saussure signe l’acte de naissance de la sémiologie et définit son champ d’étude :
\begin{lecture}{}
On peut donc concevoir une science qui étudie la vie des signes au sein de la vie sociale ; elle formerait une partie de la psychologie sociale, et par conséquent de la psychologie générale ; nous la nommerons sémiologie (du grec semeion, « signe »). Elle nous apprendrait en quoi consistent les signes, quelles lois les régissent. Puisqu’elle n’existe pas encore, on ne peut dire ce qu’elle sera ; mais elle a droit à l’existence, sa place est déterminée d’avance.
\end{lecture}
On voit dans cet extrait que pour Saussure la sémiologie n’est pas, comme le pense Peirce, une méta-science qui englobe les autres sciences. C’est plutôt une science parmi les sciences, et même une science qui appartient à une science : la psychologie sociale.

\minititle{\textcircled{\scriptsize 3}\quad La linguistique, branche et « patron » de la sémiologie}
Ferdinand de Saussure a toujours insisté sur les rapports entre linguistique et sémiologie. Par sémiologie, il entend la science sociale qui étudie les signes de manière générale. La linguistique\origcont{170}{156} n’est donc à ses yeux de Saussure qu’une branche de la sémiologie :\ednote{原文 « à ses yeux de Saussure » 如此；应删去重复成分，作 « aux yeux de Saussure » 或 « à ses yeux »。}
\begin{lecture}{}
On a discuté pour savoir si la linguistique appartenait à l’ordre des sciences naturelles ou des sciences historiques. Elle n’appartient à aucun des deux, mais à un compartiment des sciences qui, s’il n’existe pas, devrait exister sous le nom de sémiologie, c’est-à-dire science des signes ou étude de ce qui se produit lorsque l’homme essaie de signifier sa pensée au moyen d’une convention nécessaire.
\end{lecture}
Cependant, la linguistique constitue la branche la plus développée et la plus importante de la sémiologie, en raison de la complexité du langage humain, de sa « nature spéciale ». Selon lui, Saussure, la linguistique a donc vocation à devenir « le patron » de la sémiologie.
\begin{lecture}{}
La langue est une institution sociale ; mais elle se distingue par plusieurs traits des autres institutions politiques, juridiques, etc. Pour comprendre sa nature spéciale, il faut faire intervenir un nouvel ordre de faits. La langue est un système de signes exprimant des idées, et par là, comparable à l’écriture, à l’alphabet des sourds-muets, aux rites symboliques, aux formes de politesse, aux signaux militaires, etc. Elle est seulement le plus important de ces systèmes. La linguistique n’est qu’une partie de cette science générale, les lois que découvrira la sémiologie seront applicables à la linguistique, et celle-ci se trouvera ainsi rattachée à un domaine bien défini dans l’ensemble des faits humains. [...] La langue, le plus complexe et le plus répandu des systèmes d’expression, est aussi le plus caractéristique de tous ; en ce sens, la linguistique peut devenir le patron général de toute sémiologie, bien que la langue ne soit qu’un système particulier.
\end{lecture}

\Needspace{5\baselineskip}
\section{LA SÉMIOTIQUE TRIADIQUE DE PEIRCE}
Les théories du signe sont nombreuses. Nous allons commencer par présenter la sémiologie de Peirce, qui, ne le cachons pas, est un peu ardue. En comparaison, celle de Saussure, présentée aussitôt après, vous semblera un jeu d’enfant.

\subsection{Le modèle triadique}
On utilise souvent un triangle pour représenter visuellement les trois composantes d’un signe triadique : c’est ce qu’on appelle le triangle sémiotique. Si le modèle de Peirce est triadique, il s’inscrit dans une tradition fort ancienne qui remonte à Aristote.

C’est à Aristote que l’on doit le premier modèle triadique de la signification. Le texte principal en la matière est le début du \textit{De Interpretatione} :
\origpage{171}{157}
\begin{lecture}{}
La parole est un ensemble d’éléments symbolisant les états de l’âme, et l’écriture d’un ensemble d’éléments symbolisant la parole. Et, de même que les hommes n’ont pas tous le même système d’écriture, ils ne parlent pas tous de la même façon. Toutefois, ce que la parole signifie immédiatement, ce sont des états de l’âme qui, eux, sont identiques pour tous les hommes ; et ce que ces états de l’âme représentent, ce sont des choses, non moins identiques pour tout le monde.
\end{lecture}
Les « états d’âme », c’est le concept, le signifié. La parole est le signifiant. Quant aux « choses », c’est ce qu’on appelle de nos jours le référent, c’est-à-dire l’objet du monde représenté. Dans sa traduction de l’\textit{Isagoge} de Porphyre, \textbf{Boèce} retient que pour Aristote, trois facteurs, interviennent dans tout entretien et toute discussion :
\begin{itemize}
\item des choses : ce que notre esprit perçoit et que notre intellect saisi.\ednote{原文 saisi 如此；谓语动词应为 saisit。}
\item des pensées (\textit{intellectus}) : ce par quoi nous connaissons.
\item des paroles (\textit{voces}) : ce par quoi nous signifions.
\end{itemize}

\subsection{Les trois catégories phanéroscopiques}
Nous avons vu que les \textit{Catégories} d’Aristote sont au nombre de dix. Pour Peirce, seules trois catégories sont nécessaires et suffisantes pour rendre compte de toute l’expérience humaine. Pourtant, elles sont relativement difficiles à comprendre. Armons-nous de courage, et commençons par les trois « modes de l’être » (les \textit{modus essendi} des modistes) de Peirce :
\begin{lecture}{}
Mon opinion est qu’il y a trois modes d’être. Je soutiens que nous pouvons les observer directement dans les éléments de tout ce qui est à n’importe quel moment présent à l’esprit d’une façon ou d’une autre. Ce sont l’être de la possibilité qualitative positive, l’être du fait actuel (c’est-à-dire « en acte »), et l’être de la loi qui gouvernera les faits dans le futur.
\end{lecture}
\minititle{\textcircled{\scriptsize 1}\quad La notion de \textit{phaneron}}
Dans la définition ci-dessus, « tout ce qui est présent à l’esprit, que cela corresponde à une chose réelle ou pas » correspond à ce que Peirce appelle « \textit{phanéron} », mot grec signifiant « ce qui est visible ». Sa sémiologie est parfois qualifiée de phanéroscopie, c’est-à-dire l’étude des phanérons.\ednote{需区分两门在 Peirce 体系中相关却不同的学科：phaneroscopy（现象学）研究经验的一般范畴；逻辑／符号学则属规范科学。不能把两者无条件作为同义词。见 Peirce 的 1903 年科学分类，\ecite{37}。} On pourrait aussi parler de phénoménologie, bien que ce terme soit habituellement associé au philosophe et logicien \textbf{Edmund Husserl} (1859-1938). Peirce définit ainsi le phanéron :
\begin{lecture}{}
Je propose d’utiliser le mot Phanéron comme un nom propre pour dénoter le contenu total d’une conscience [...] la somme de tout ce que nous avons à l’esprit, de quelque manière que ce soit, sans regarder sa valeur cognitive. Ceci est assez vague : mais c’est volontaire, je soulignerai seulement que je ne limite pas la référence à un état de conscience instantané ; car la clause “de quelque manière que ce soit” comprend la mémoire et toute cognition habituelle.
\end{lecture}

\origpage{172}{158}
Pour Peirce, tout phanéron peut être analysé avec trois catégories : la priméité (\textit{firstness}), la secondéité (\textit{secondness}) et tiercéité (\textit{thirdness}).

\minititle{1. La Priméité}
La Priméité (\textit{firstness}), c’est l’essence, au sens philosophique, de tout ce qui est. C’est la conception de l’être indépendamment de toute autre chose réelle et indépendamment de notre pensée. C’est le monde des Possibles, des potentialités abstraites, des qualités générales non matérialisées mais pouvant l’être, comme les couleurs ou les formes. Comme exemple de \textit{premier}, Peirce cite une qualité pure ou un sentiment : la « rougéité », qualité d’être rouge, ou le sentiment de cette « rougéité » avant qu’elle se matérialise, qu’elle s’incarne dans un objet. Le « premier » correspond à la vie émotionnelle. À la notion de priméité sont attachées les idées de commencement, d’indifférenciation. Dans la philosophie grecque, elle correspond à l’\textit{oussia} (l’essence) d’Aristote.

\minititle{2. La Secondéité}
Si la Priméité est le monde des Possibles et la catégorie des Essences, la secondéité (\textit{secondness}) est l’univers des Existants, la catégorie de l’Existence de tout ce qui est. C’est la catégorie des faits actuels, de l’individuel, de l’expérience, de l’action, de la réaction, de la vie pratique. Par exemple, la « rougéité » (qualité d’être rouge, donc « premier ») s’incarne dans un tableau, objet individuel, réel. Dans la terminologie de Peirce, le tableau est un « second ». Dans la philosophie grecque, la Secondéité correspond à la \textit{tyché} (l’accident) d’Aristote.

\minititle{3. La Tiercéité}
La tiercéité (\textit{thirdness}) est la médiation qui met en relation le monde des Possibles et le monde des Existants. Peirce nous en propose quelques exemples : une route entre deux points, un messager, le moyen terme d’un syllogisme, un interprète. Le « troisième » est de l’ordre des conventions et des lois. Le langage, fait des règles, est un « troisième ».\ednote{原文 fait des règles 如此；按语意应为 fait de règles（由规则构成）。} La tiercéité est la catégorie du langage, et donc de la pensée, de la représentation, et du processus sémiotique : elle correspond à la vie intellectuelle. Dans la philosophie grecque, la Secondéité correspond au \textit{logos}.\ednote{此处 Secondéité 与本小节论述及前一句不符，应为 Tiercéité。与 Aristote 术语的对应是作者的解释性类比，不宜当作三套范畴严格同一。}

Notez qu’une « loi » ne se manifeste qu’à travers des « faits » qui l’appliquent, donc dans la secondéité. Et ces faits eux-mêmes sont les incarnations (actualisations) de qualités, donc de la priméité. On peut résumer cette dynamique de la façon suivante : la pensée (tiercéité) saisit les existants (secondéités) comme des possibles (priméités) réalisés.

\subsection{La tripartition du signe}
Au cours de sa vie longue et prolifique, Charles Sanders Peirce a donné près de quatre-vingts définitions du concept de signe, ce qui illustre la difficulté d’utiliser des signes pour définir d’autres signes. Nous allons en voir trois qui ont l’avantage de contenir toutes les autres.

\minititle{\textcircled{\scriptsize 1}\quad Définitions du signe}
Commençons par une définition de 1868 :
\origpage{173}{159}
\begin{lecture}{}
Un signe a, comme tel, trois références : premièrement il est un signe pour quelque pensée qui l’interprète ; deuxièmement, il est un signe qui tient lieu de quelque objet auquel il est équivalent dans cette pensée ; troisièmement il est un signe, sous quelque rapport ou qualité, qui le met en connexion avec cet objet.
\end{lecture}
Dans cette définition, le signe est défini implicitement comme une triade (« trois références »). On retrouve aussi ce rapport triadique (« trois choses ») dans la définition de 1873 :
\begin{lecture}{}
Un signe est quelque chose qui représente une autre chose pour un esprit. Pour son existence comme tel trois choses sont requises. En premier lieu, il doit avoir des caractères qui nous permettront de le distinguer des autres objets. En second lieu, il doit être affecté d’une façon ou d’une autre par l’objet qui est signifié. En troisième lieu chaque représentation s’adresse à un esprit. C’est seulement pour autant qu’elle fait ceci qu’elle est une représentation.
\end{lecture}
Enfin, dans ses \textit{Elements of Logic} (1903), il reprend et complète l’adage latin qui définit le signe comme « \textit{aliquid stat pro aliquo} », c’est-à-dire une chose qui tient lieu d’une autre.
\begin{lecture}{}
Un signe, ou representamen, est quelque chose qui tient lieu, pour quelqu’un de quelque chose, sous quelque rapport ou à quelque titre. Il s’adresse à quelqu’un, c’est à dire qu’il crée, dans l’esprit de cette personne un signe plus développé. Ce signe qu’il crée, je l’appelle l’interprétant du premier.
\end{lecture}
On voit dans cette troisième définition que le processus sémiotique (appelé « sémiose ») nécessite la coopération des trois instances que sont le signe (ce qui représente), l’objet (ce qui est représenté) et l’interprétant, qui produit leur relation. Le rapport de sémiose, quant à lui, désigne une action ou une influence qui est ou qui suppose la coopération de trois sujets, tels que le signe, son objet et son interprétant.

Prenons comme exemple un panneau routier (un stop) sur la route. Le panneau de fer est le \textit{representamen}. Il est mis sur le bord de la route pour indiquer un ordre d’arrêt : c’est l’\textit{objet} du signe. Il crée dans l’esprit de celui qui le perçoit un autre signe m’indiquant ce qu’il doit faire : c’est l’\textit{interprétant} du signe. Mais selon que celui qui perçoit le signe est un automobiliste ou un piéton, l’\textit{interprétant} ne fonctionnera pas de la même façon : la sémiosis est liée au contexte. Examinons ces trois éléments plus en détail.

\minititle{\textcircled{\scriptsize 2}\quad Le \textit{representamen} (« signe »)}
Le \textit{representamen} est le signe matériel qui représente (dénote) une autre chose, à savoir un objet de pensée. Avant d’être interprété, il n’est qu’une pure possibilité de signifier : c’est donc, selon les catégories phanéroscopiques, un \textit{premier}.\ednote{此处“物质符号”把范围说窄了。Peirce 的符号可以是性质、个别事物，也可以是一般规律；representamen 在三元关系中充当第一项，不等于它本身必然只是物质实体或 qualisign。下文的三分表也应据此理解。见\ecite{37}。}

\origpage{174}{160}
\minititle{\textcircled{\scriptsize 3}\quad L’objet}
L’objet est ce que le signe représente. Pour Peirce, l’objet renvoie à « tout ce qui vient à la pensée ou à l’esprit dans le monde ordinaire ». Comme il représente ce qui existe et ce dont on parle, l’objet est donc un \textit{second}. Par ailleurs, on vient de le voir un peu plus haut, la notion d’objet est proche de celle de contexte :
\begin{lecture}{}
Il faut distinguer deux notions d’objets, c’est-à-dire deux sortes de contexte : le contexte intrinsèque, tel que le signe le représente (objet immédiat) et le contexte extrinsèque, c’est-à-dire pris en lui-même, extérieurement au signe-chose (objet dynamique).
\end{lecture}
Peirce distingue donc :
\begin{itemize}
\item l’objet immédiat : l’objet tel que le signe le représente
\item l’objet dynamique : l’objet tel qu’il est dans la réalité
\end{itemize}
Le signe ne peut que représenter l’objet : il ne peut pas le faire connaître.
\begin{lecture}{}
La réalité est ce mode d’être en vertu duquel la chose réelle est comme elle est, indépendamment de ce qu’un esprit ou une collection déterminée d’esprits peuvent se représenter qu’elle est.
\end{lecture}
Ce que le signe ne peut pas directement exprimer mais ne peut qu’indiquer, le récepteur doit l’interpréter grâce à son expérience.

\minititle{\textcircled{\scriptsize 4}\quad L’« interprétant »}
Exprimé de façon simple, l’interprétant est la pensée (la représentation mentale) qui permet de renvoyer le \textit{representamen} à son objet. Peirce le définit comme « l’effet du signe (c’est-à-dire du representamen sur une personne » :\ednote{原文括号未闭合，宜在 representamen 后补右括号，作 « l’effet du signe (c’est-à-dire du representamen) sur une personne »。此为标点拟校，不据此宣称已核定引文的原始出处。} autrement dit, l’interprétant est la valeur que prend le representamen dès qu’il est perçu (« effet ») par un sujet. C’est un \textit{troisième} qui dynamise la relation de signification. Par exemple, si je parle d’un livre, le mot « livre » est le \textit{représentamen}, l’\textit{objet} est ce qui est désigné par ce mot, et l’\textit{interprétant} est la définition communément admise de ce mot, à savoir le concept de livre. Peirce distingue :
\begin{itemize}
\item L’interprétant immédiat, c’est-à-dire le sens commun, sens minimal, susceptible de venir spontanément à l’esprit de n’importe quel sujet qui perçoit le signe.
\item L’interprétant dynamique, qui est le sens particulier formé dans l’esprit d’un récepteur particulier.
\item L’interprétant final, sens sur lequel tous les récepteurs peuvent s’accorder. On peut dire que c’est le sens « correct » ou « autorisé ».
\end{itemize}
\minititle{\textcircled{\scriptsize 5}\quad Illustration}
Dans \textit{The American Journal Of Semiotics} (1987), \textbf{Nathan Houser}, professeur de philosophie à\origcont{175}{161} l’université d’Indianapolis et grand connaisseur de Peirce propose de comprendre ce modèle avec un exemple médical :

Un patient se présente chez le médecin avec de la fièvre et la gorge enflammée, symptômes qui constituent le signe.

Le médecin, connaissant un certain nombre de maladies qui provoquent ces symptômes, formule un diagnostic : par exemple, « c’est un rhume ». Le rhume, maladie la plus facilement associée à ces symptômes, constitue l’\textit{objet immédiat}, alors que le diagnostic lui-même constitue l’\textit{interprétant immédiat}.

Le médecin fait alors une ordonnance et un pronostic (« Vous irez mieux dans une semaine. ») qui constituent l’\textit{interprétant dynamique}. Dans ce cas, l’\textit{objet dynamique} est la maladie qui a véritablement causé les symptômes (qu’il s’agisse de celle diagnostiquée par le médecin ou d’une autre présentant les mêmes symptômes), tandis que l’\textit{interprétant final} serait le diagnostic correct.

\minititle{\textcircled{\scriptsize 6}\quad La sémiosis : le processus de signification}
Selon Peirce, la sémiosis est un processus de signification qui se déroule dans l’esprit de l’interprète :
\begin{lecture}{}
Il débute avec la perception du signe et se termine avec la présence à son esprit de l’objet du signe. C’est un processus inférentiel.
\end{lecture}
Le \textit{representamen}, pris en considération par un interprète, a le pouvoir de déclencher un interprétant, qui est lui-même un signe (un representamen) et renvoie, par l’intermédiaire d’un autre interprétant, au même objet que le premier representamen, permettant ainsi à ce premier de renvoyer à l’objet. Et ainsi de suite, ce qui semble être un processus infini. Prenons l’exemple de la définition d’un mot dans le dictionnaire : la définition de ce mot est un \textit{interprétant}, puisqu’elle renvoie à l’objet (ce que représente ce mot) et permet donc au \textit{representamen} (le mot) de renvoyer à cet objet. Mais la définition elle-même, pour être comprise, nécessite une série d’autres interprétants (d’autres définitions)... Ainsi, le processus sémiotique est, théoriquement, illimité.

Le modèle de Peirce inscrit le processus de sémiosis dans un cadre pragmatique et communicatif. La distinction entre l’immédiat et le dynamique permet de séparer ce qui est reconnu, décodé de façon mécanique, et l’interprétation effective par un récepteur. Par exemple, pour reprendre l’exemple ci-dessus, la définition d’un mot dans un dictionnaire (interprétant \textit{immédiat}) et le sens de ce mot en contexte (interprétant \textit{dynamique}).

L’interprète, lieu de ce processus de sémiosis, réagit au signe en fonction de son expérience personnelle (interprétant dynamique). Il est donc possible pour un même \textit{representamen} qu’un sujet donné ait une interprétation très différente de celle d’un autre, dans la mesure où les interprétants dont ils sont porteurs peuvent être différents. Les interprétants « émotionnels », par exemple, relevant de la priméité, peuvent avoir un effet sur le jugement. Les interprétants « logiques », peuvent aussi varier d’une personne à l’autre.

\origpage{176}{162}
Par « signe », Peirce entend à la fois une potentialité signifiante et un acte interprétatif particulier. On peut ainsi déjà saisir la nuance fondamentale entre la dénotation (le sens qui est « dans » le signe), et la connotation, sens « hors » du signe, imposé par la culture et ne dépendant donc pas d’une interprétation purement personnelle\origfoot{2}{Nous y reviendrons dans le chapitre 11 consacré à la sémantique.}.

\subsection{La trichotomie du signe}
Le signe, on vient de le voir, est une triade, un processus de relations entre un representamen, un objet et un interprétant. Pour caractériser le signe dans son individualité, il est nécessaire de le trichotomiser, c’est-à-dire d’envisager chacune de ses trois composantes du point de vue de la catégorie phanéroscopique dont elle relève. Peirce a ainsi identifié neuf sous-signes qui caractérisent ce qu’il appelle les divers « moments »et rapports du signe :\ednote{应区分符号的三项关系与三种分类依据。表中是三组三分：符号自身、符号与对象的关系、符号与解释项的关系；不是把一个符号拆成九个独立“子符号”。1903 年分类在相容限制下给出十类复合符号。见\ecite{37}。}
\begin{center}\small
\begin{tabularx}{\linewidth}{@{}l*{3}{>{\centering\arraybackslash}X}@{}}
\toprule
 & \textbf{Priméité} & \textbf{Secondéité} & \textbf{Tiercéité}\\\midrule
Representamen & Qualisigne & Sinsigne & Légisigne\\
Objet & Icône & Indice & Symbole\\
Interprétant & Rhème & Dicisigne & Argument\\\bottomrule
\end{tabularx}
\end{center}
\minititle{\textcircled{\scriptsize 1}\quad La trichotomie du representamen}
Chacun des trois termes du processus sémiotique se subdivise selon les trois catégories, et Peirce distingue :

\minititle{1. Le qualisigne}
Un \textit{qualisigne} (priméité) est une qualité qui fonctionne comme signe, comme par exemple la couleur rouge.

\minititle{2. Le sinsigne}
Un \textit{sinsigne} (secondéité) est une chose déterminée qui fonctionne comme signe, comme un feu rouge. Peirce en donne la définition suivante : « Un sinsigne est une chose ou un événement existant réel qui est un signe ».

\minititle{3. Le légisigne}
Un \textit{légisigne} (tiercéité) : une loi qui est un signe, c’est-à-dire un signe conventionnel. Par exemple : l’interdiction de franchir un feu rouge. Un mot de passe, un insigne militaire, un billet de banque, les signaux du code de la route, les systèmes d’écriture sont des exemples de légisignes basés sur la loi, la convention. La langue étant une convention, les mots de la langue sont donc des légisignes. Pour Peirce, quand nous disons que « le » est un « mot », que « un » est un autre « mot », alors un « mot » est un « légisigne ».

\origpage{177}{163}
Ces légisignes ne peuvent agir qu’en se matérialisant dans des sinsignes qui constituent des « répliques ». Par exemple, l’article « le » est un légisigne, mais il ne peut être matérialisé (vu ou entendu) que si on le prononce ou ‘écrit.\ednote{原文 « ou ‘écrit » 如此；应为 « ou l’écrit »。} La voix et l’écriture, matérialisent le \textit{légisigne} « le » dans un \textit{sinsigne}. Mais ce n’est pas tout : le \textit{sinsigne} « le », selon qu’il est écrit ou prononcé, comprend également des \textit{qualisignes} : l’intonation, le timbre de la voix, ou la forme des lettres à l’écrit.

\minititle{\textcircled{\scriptsize 2}\quad La trichotomie de l’objet}
Un \textit{representamen} peut renvoyer à son objet selon la priméité, la secondéité ou la tiercéité, c’est-à-dire dans un rapport de similarité, de contiguïté, et de loi. Suivant cette trichotomie, le signe est appelé respectivement : icone, indice et symbole.

\minititle{1. L’icone}
Exprimé en termes simples, l’icone (du genre masculin chez Peirce, et sans accent circonflexe) est un signe qui renvoie à l’objet signifié au moyen d’une similarité ou d’une ressemblance avec celui-ci :
\begin{lecture}{}
Un icone est un signe qui fait référence à l’Objet qu’il dénote simplement en vertu de ses caractères propres, lesquels il possède, qu’un tel Objet existe réellement ou non. [...] N’importe quoi, que ce soit une qualité, un existant individuel, ou une loi, est un icone de n’importe quoi, dans la mesure où il ressemble à cette chose et en est utilisé comme le signe.
\end{lecture}
Par exemple sur une photographie ou sur un tableau, le portrait (l’icone) renvoie à l’objet. De même, évoquer par exemple une couleur au moyen d’un objet (« vert émeraude », « rose bonbon », « vert pomme ») est un processus iconique. Une maquette ou la miniature d’un bâtiment sont aussi des icones. Mais attention : si le dessin d’un verre est bien un icone, cet icone placé sur une caisse (lors d’un déménagement par exemple) devient un « légisigne », qui, par convention, signifie « fragile ».\ednote{icone 与 légisigne 属于不同分类维度，并不互斥；一个一般性的图示可以同时是 iconic legisign。不能据“约定使用”便推断其像似关系消失。见\ecite{37}，1903 年十类符号中的第五类。}

\minititle{2. L’indice}
L’objet, envisagé sous l’angle de la secondéité, est appelé indice.\ednote{这里不是对象本身变为 indice，而是符号因与对象有现实联系而被归为 indice。下引 Peirce 的定义也以“Un indice est un signe”开头。见\ecite{37}。} En voici la définition qu’en donne Peirce :
\begin{lecture}{}
Un indice est un signe qui fait référence à l’Objet qu’il dénote en vertu du fait qu’il est réellement affecté par cet Objet. Dans la mesure où l’Indice est affecté par l’Objet, il a nécessairement certaines qualités en commun avec cet Objet, et c’est sous ce rapport qu’il réfère à l’Objet. Il implique, par conséquent, une certaine relation iconique à l’Objet, mais un icone d’un genre particulier ; et ce n’est pas la simple ressemblance à son Objet, même sous ces rapports, qui en font un signe mais les modifications réelles qu’il subit de la part de l’Objet.
\end{lecture}
Cette définition étant assez obscure, nous pouvons faire référence à celle du sémiologue argentin \textbf{Luis Jorge Prieto} (1926-1996) tirée de son \textit{Essai de sémiologie}(1975) : « Fait immédiatement\origcont{178}{164} perceptible qui nous fait connaître quelque chose à propos d’un autre fait qui ne l’est pas. » L’étymologie est aussi fort utile : « indice » est dérivé du latin \textit{index}, et l’index est le nom du doigt que l’on utilise pour montrer : l’indice est donc un signe révélateur. Si, comme tout signe, il produit du sens, il le fait à l’insu du producteur de sens. Ceci est une caractéristique essentielle de l’indice : c’est un signe non intentionnel.\ednote{“非意图性”不是 Peirce 的 index 的必要条件；他也将有意指点的手指、指示词等列入指示性符号。这里把 Prieto 式 signal／indice 区分直接移用于 Peirce，宜分别说明其理论框架。见\ecite{37}。} Dans l’exemple donné par Nathan Houser un peu plus haut, le mal de gorge est l’indice (le symptôme) de la maladie. Des nuages noirs à l’horizon sont l’indice de la pluie à venir, la fumée est l’indice du feu (« il n’y a pas de fumée sans feu »). Le mouvement d’une girouette est causé par la direction du vent : elle en est l’indice. Un coup frappé à la porte est l’indice d’une visite.

\minititle{3. Le Symbole}
Le troisième et dernier type de representamen dans la trichotomie de Peirce est le symbole, qu’il définit ainsi :
\begin{lecture}{}
Un signe qui se réfère à l’Objet qu’il dénote en vertu d’une loi, habituellement une association générale d’idées, qui provoque le fait que le Symbole est interprété comme référant à l’Objet.
\end{lecture}
Le symbole étant une « tiercéité », on ne sera pas étonné de voir Peirce le définir comme un signe renvoyant à son objet en vertu d’une « loi ». Ainsi donc, les légisignes rencontrés plus haut (ticket d’entrée à un spectacle, un billet de banque, les mots de la langue, etc.) sont tous des symboles. La « règle » symbolique peut avoir été formulée \textit{a priori}, par convention, ou s’être constituée \textit{a posteriori} par habitude culturelle. Pour \textbf{Luis Jorge Prieto}, « un symbole est «la notion d’un rapport dans une culture donnée entre deux éléments ». Ici encore, l’étymologie du mot « symbole » est éclairante : à l’origine, un \textit{symbolon} (en grec, « joindre, relier ») était un objet partagé en deux et dont la possession de chacune des deux parties permettait à deux individus différents de se reconnaître. Charles Peirce explique :\ednote{底本在“Charles Peirce explique :”之后未另列可辨识的引文，下一段以“D’ailleurs que”起。疑有脱漏或引文标识缺失；未取得可确证的对勘材料，不补造文字。}

D’ailleurs que lorsque deux signes sont « reliés » par leur signifiés, on parle de symbolisation : dans cette optique, un mot de passe, ou un code sont des exemples de symbole. La langue étant un code, on comprend alors qu’elle soit symbolique, puisque le signe linguistique est lui-même un type particulier de symbole.

Les symboles sont si nombreux que nous ne pouvons n’en citer que quelques-uns.\ednote{原文 « ne pouvons n’en citer que » 如此；两种规范表达分别是 « nous ne pouvons en citer que quelques-uns » 和 « nous n’en citerons que quelques-uns »。} La balance et le glaive sont deux symboles différents de la justice, reliés l’un et l’autre à des valeurs culturelles très fortes : l’équité pour la balance, et la force pour le glaive. La colombe est le symbole de la paix, le lion est le symbole du courage, la croix latine est le symbole du christianisme, le sceptre et la couronne sont les symboles de la royauté. Contrairement à ce que nous avons vu à propos de l’indice, la fonction première du symbole est la fonction sémiotique, c’est-à-dire signifier.

\minititle{4. Le signal}
Afin d’être complet, et avant de conclure sur la trichotomie du signe, il convient définir le signal,\ednote{原文 « il convient définir » 如此；应为 « il convient de définir »。}\origcont{179}{165} bien qu’il n’appartienne pas à la trichotomie des signes de Peirce. Le but est de montrer en quoi il diffère des autres types de signes. La notion de signal ne nous est pas étrangère : dans le chapitre 2, nous avons étudié les différents types de signaux utilisés dans le monde animal, et nous avons appris avec Benveniste que le mode de communication des abeilles est « code de signaux ». Mais qu’est-ce qu’un signal ? Pour Prieto, c’est « un fait qui a été produit artificiellement pour servir d’indice ». Dit plus simplement, c’est une information donnée par un émetteur qui aide le récepteur à prendre une décision. En d’autres termes, c’est la représentation physique d’une information à transmettre. La sonnerie qui résonne à la fin du cours est un signal attendu et apprécié de tous les étudiants : c’est le signal que le cours est fini. Un coup de pistolet tiré au début d’une course indique le début de la course. On le voit : l’intention de communication, de signification est essentielle, et c’est elle qui différencie le signal de l’indice. La fumée, on l’a vu un peu plus haut, est l’indice du feu. Mais lorsque le roi You de Zhou (周幽王) fait allumer un feu pour faire sourire la belle Bao Si (褒姒), le feu devient un signal.

\minititle{\textcircled{\scriptsize 3}\quad La trichotomie de l’interprétant}
Suivant la trichotomie de l’interprétant, le signe est appelé respectivement rhème (qui décrit une possibilité d’essence), dicisigne (qui décrit une existence) ou argument (ou raisonnement).

\minititle{1. Le rhème}
L’interprétant rhématique a une structure de priméité. Pour opérer la relation du representamen à l’objet, il ne fait appel qu’aux qualités du representamen. Ces qualités sont les qualités de toute une classe d’objets possibles. Par conséquent, le rhème, qui correspond à une variable dans une fonction propositionnelle, n’est ni vrai ni faux. Il fonctionne comme un blanc dans une formule, un vide à remplir dans un exercice à trous, de type : « ... est rouge ». Le portrait-robot d’une personne, sans autre indication, représente toute une classe d’objets possibles : les personnes ressemblant à ce portrait.

\minititle{2. Le dicisigne}
Le dicisigne est un signe interprété au niveau de la secondéité. Il fonctionne comme une équation ou une proposition logique qui met en relation des constantes (un sujet et prédicat), et qui peut être vraie ou fausse. Par exemple, le portrait d’une personne avec l’indication du nom de cette personne est un \textit{dicisigne}. Un dicisigne ne fournit pas de raison de sa vérité ou de sa fausseté, contrairement à un argument qui aboutit à une conclusion en suivant un processus rationnel.

\minititle{3. L’argument}
L’argument interprète un signe au niveau de la tiercéité et formule la règle qui relie le representamen et son objet. On distingue trois types d’arguments selon la nature de la règle qui relie le representamen à son objet. Ces trois arguments, que nous avons déjà rencontrés dans le chapitre 4, sont la déduction, l’induction et l’abduction, « la seule espèce de raisonnement susceptible d’introduire des idées nouvelles ».

\origpage{180}{166}
\subsection{Charles W. Morris : le successeur}
Si l’on reconnaît en Peirce un précurseur de Karl Popper, il a aussi directement inspiré les philosophes pragmatistes américains \textbf{William James} (1842-1910) et \textbf{John Dewey} (1859-1952). En logique et philosophie des mathématiques, son influence a aussi été marquante sur \textbf{Willard Van Orman Quine} (1908-2000) et \textbf{Hilary Putnam} (1926-2019).\ednote{原文 1926-2019 如此；Hilary Putnam 于 2016 年 3 月 13 日去世，生卒年应为 1926-2016。见 Harvard 纪念资料，\ecite{40}。} En sémiotique, l’œuvre de Charles W. Morris est considérée comme un commentaire de celle de Peirce.

\textbf{Charles W. Morris} (1903-1979), est un sémioticien et philosophe américain. Dans son ouvrage \textit{Foundations of the Theory of Signs} (1938), sur la base de la théorie de Peirce, et en se référant à « une tradition qui remonte aux Grecs », il distingue d’abord trois composants dans la sémiosis :
\begin{itemize}
\item le \textit{sign-vehicle} (vecteur du signe) : ce qui agit comme signe
\item le \textit{designatum} (le designé) : ce à quoi le signe réfère
\item le \textit{interpretant} : effet produit sur un certain interprète
\end{itemize}
Dans une version révisée de son modèle, Morris a intégré l’interprète : « l’interprète peut être inclus en qualité de quatrième facteur ».\ednote{待核：本句把解释者作为“第四因素”的说明归于一个修订版本，却未列该版年份及页码。本版尚未取得可与此句对应的修订版原文，不据此确认其先后关系。} Pour rappel, l’interprète est la personne pour qui le signe a fonction de signe. On passe donc d’un modèle triadique à un modèle tétradique. Comme chez Peirce, l’interprétant est l’expression d’un rapport dialectique entre le monde réel, un « déjà-là » codifié par des rapports institutionnalisés, intériorisé par les individus (par l’inculcation pédagogique qu’exerce la société institutrice) et les individus eux-mêmes qui les actualisent dans leurs déterminations concrètes par un acte interprétatif. L’interprète est un individu singulier qui est le lieu, le support de cet acte interprétatif. En tant qu’être social il est porteur des valeurs dominantes de la société et donc des significations collectivement accordées aux choses.

\section{LE SIGNE DYADIQUE DE SAUSSURE}
Après un long passage outre-Atlantique revenons en Europe pour nous recentrer sur la linguistique saussurienne.

\subsection{Caractéristiques du signe linguistique}
On l’a vu dans les pages précédentes : le signe linguistique a reçu de nombreuses définitions, notamment constitutives. Un signe se reconnaît à la présence de termes particuliers et de relations particulières entre ces termes. Pour Saussure, le signe est dual.

\minititle{\textcircled{\scriptsize 1}\quad Dualité du signe linguistique : signifiant et signifié}
Ferdinand de Saussure a proposé au début du 20\textsuperscript{e} siècle un signe dyadique, composé d’un signifiant et d’un signifié, souvent représenté par un cercle dont la moitié inférieure est le signifiant et la moitié supérieure le signifié. Selon lui, le signe linguistique unit, « non pas un nom et une chose, mais un concept et une image acoustique ». L’image acoustique (le signifiant) n’est pas le son matériel mais l’empreinte psychique de ce son dans notre esprit. Le concept (le signifié) contient
\origcont{181}{167} les traits (éléments) distinctifs qui caractérisent ce signe par rapport aux traits d’autres signes de la langue. Le signe linguistique selon Saussure se définit donc comme une entité psychique à deux visages. Par exemple, le mot français « arbre » est un signe linguistique associant la forme sonore ({\BBSixGreek /aʁbʁ/}) au concept d’arbre.
\begin{lecture}{}
Pour certaines personnes la langue, ramenée à son principe essentiel, est une nomenclature, c’est-à-dire une liste de termes correspondant à autant de choses. Cette conception est critiquable à bien des égards. [...] Cependant cette vue simpliste peut nous approcher de la vérité, en nous montrant que l’unité linguistique est une chose double, faite du rapprochement de deux termes. Le signe linguistique unit non une chose et un nom mais un concept et une image acoustique. Cette dernière n’est pas le son matériel, chose purement physique, mais l’empreinte psychique de ce son, la représentation que nous en donne le témoignage de nos sens ; elle est sensorielle, et s’il nous arrive de l’appeler matérielle, c’est seulement dans ce sens et par opposition à l’autre terme de l’association, le concept, généralement plus abstrait. [...]

Le signe linguistique est donc une entité psychique à deux faces. [...]Nous proposons de conserver le mot signe pour désigner le total, et de remplacer concept et image acoustique respectivement par signifié et signifiant ; ces derniers termes ont l’avantage de marquer l’opposition qui les sépare soit entre eux, soit du total dont ils font partie. Quant à signe, si nous nous en contentons, c’est que nous ne savons par quoi le remplacer, la langue « usuelle » n’en suggérant aucun autre. Le signe linguistique ainsi défini possède deux caractères primordiaux.
\end{lecture}
Pour illustrer de façon imagée la dualité du signe linguistique, Saussure utilise la comparaison avec une feuille de papier :
\begin{lecture}{}
La langue est encore comparable à une feuille de papier : la pensée est le recto et le son le verso ; on ne peut découper le recto sans découper en même temps le verso ; de même dans la langue, on ne saurait isoler ni le son de la pensée, ni la pensée du son ; on n’y arriverait que par une abstraction dont le résultat serait de faire de la psychologie pure ou de la phonologie pure.
\end{lecture}
Au-delà de l’efficacité de l’image utilisée, le court extrait ci-dessus nous rappelle que la théorie du signe de Saussure est basée sur des processus phonologiques mais aussi psychiques. Pour cette raison, la linguistique saussurienne est généralement qualifiée de mentaliste.

\minititle{\textcircled{\scriptsize 2}\quad Arbitraire du signe linguistique}
La notion d’arbitraire du signe linguistique, habituellement associée à Saussure, n’est pas récente : Bien avant cela, l’arbitraire du signe était évoqué dès l’antiquité, en Inde et en Grèce. Un dialogue de Platon portant sur la rectitude des noms, le \textit{Cratyle} ; a joué un rôle considérable et durable dans la philosophie du langage en Occident.\ednote{原文在 Cratyle 后用分号，隔断主语与谓语；应为逗号。}

\origpage{182}{168}
\minititle{1. L’arbitraire du signe chez Platon : le \textit{Cratyle}}
En Grèce, l’adoption de l’alphabet d’une langue sémitique (le phénicien) après l’abandon du mycénien, ainsi que la naissance de la rhétorique, ont été à l’origine de l’intérêt des philosophes pour le langage. Ainsi, à la fin du 6\textsuperscript{e} siècle av. J.-C., le philosophe présocratique grec \textbf{Héraclite d’Éphèse} distinguait la pensée (\textit{logos}), l’énoncé (\textit{epos}) et la réalité (\textit{ergon}), et considérait que le lien entre ces trois éléments était réalisé par le \textit{logos}, principe divin, unique. Au 4\textsuperscript{e} siècle avant J.-C., \textbf{Démocrite}, philosophe matérialiste et père de l’atomisme, s’est opposé à l’origine divine du langage : pour lui, le langage est purement conventionnel. Dans son \textit{Cratyle}, Platon expose deux thèses opposées sur la nature des mots.

\minititle{1) La thèse d’Hermogène}
Pour Hermogène, partisan de l’arbitraire du signe, il n’y a qu’un lien abstrait et extrinsèque entre la chose et le nom, et ce lien est établi par convention. Les mots sont attribués aux choses par une sorte de « législateur » :
\begin{lecture}{}
Pour moi, Socrate, après en avoir souvent raisonné avec Cratyle et avec beaucoup d’autres, je ne saurais me persuader que la propriété du nom réside ailleurs que dans la convention et le consentement des hommes. Je pense que le vrai nom d’un objet est celui qu’on lui impose ; que si à ce nom on en substitue un autre, ce dernier n’est pas moins propre que n’était le précédent : de même que si nous venons à changer les noms de nos esclaves, les nouveaux qu’il nous plaît de leur donner ne valent pas moins que les anciens. Je pense qu’il n’y a pas de nom qui soit naturellement propre à une chose plutôt qu’à une autre, et que c’est la loi et l’usage qui les ont tous établis et consacrés.
\end{lecture}
Derrière la pensée d’Hermogène, on retrouve la thèse du sophiste grec \textbf{Protagoras d’Abdère} (vers 490-420 avant J.-C.) pour qui « l’homme est la mesure de toute chose ». Appliquée au langage, cette thèse affirme que c’est l’homme qui donne un sens à toute chose.

\minititle{2) La thèse de Cratyle}
Cratyle, disciple d’Héraclite, défend au contraire la théorie d’une relation motivée, naturelle, entre les mots et les choses : pour lui, les mots sont une peinture des choses. Ils ressemblent à ce qu’ils signifient -- ce sont des symboles et les noms sont justes par nature :
\begin{lecture}{}
Il existe une dénomination naturelle pour chacun des êtres [...] Un nom n’est pas l’appellation que certains donnent à l’objet après accord, en le désignant par une parcelle de leur langage, mais, il existe naturellement, et pour les Grecs et pour les Barbares, une juste façon de dénommer qui est la même pour tous.
\end{lecture}

\origpage{183}{169}
\minititle{3) La synthèse de Socrate}
Contre Hermogène, Socrate établit d’abord que les mots sont des instruments qui servent à nommer la réalité : ils ont donc un lien avec elles.\ednote{原文 elles 如此；照应单数 la réalité 应为 elle。} Le nom est l’instrument qui permet de nommer, et c’est un législateur qui établit les noms et compose, à partir de syllabes, le nom qui correspond à une chose. On peut juger de la justesse (de la rectitude) du nom grâce à l’étymologie, qui permet de retrouver dans les noms le \textit{logos} héraclitéen. Par exemple, les hommes, admirant les astres du ciel toujours en train de courir (\textit{thein}, en grec), appelèrent les dieux \textit{theos}. Ainsi, de mot en mot, l’étymologie permet de remonter aux noms primitifs, et les noms, par les lettres et les syllabes, imitent la nature d’un objet pour la nommer : la lettre {\BBSixGreek ρ} (\textit{rhô}) suggère l’expression du mouvement, le {\BBSixGreek δ} (\textit{delta}) ou le {\BBSixGreek τ} (\textit{tau}) expriment l’enchaînement ou l’arrêt, tandis d’autres lettres imitent le bruit du vent ou l’écoulement de l’eau.\ednote{原文 tandis d’autres 如此；应为 tandis que d’autres。本段复述对话中的词源说明，不应把它们当作现代历史语言学已确认的词源。} C’est la thèse de Cratyle. En plus de nommer la réalité, les mots sont aussi des images de réalité, un peu comme des peintures, qui sont également des imitations d’un autre genre.

Mais une image n’imite jamais parfaitement une chose, sinon ce n’est plus une image mais une copie. Si le nom de Cratyle imitait parfaitement Cratyle, il n’y aurait plus un mais deux Cratyle. Le mot, comme la peinture, possède donc des imperfections nécessaires pour ne pas redoubler les choses d’une autre réalité faites de mots.\ednote{原文 faites 如此；与单数 réalité 一致应为 faite。} Contre Cratyle, Socrate établit que le nom ne doit pas être exactement la chose mais simplement désigner les caractéristiques d’une chose ou la chose en soi. Devant Cratyle, qui peine à l’admettre, Socrate montre la part de convention dans les noms : c’est l’usage qui parfois se substitue à la ressemblance pour désigner une chose. Dans son \textit{Cours de linguistique générale}, nous allons voir Saussure s’est rangé du côté d’Hermogène.\ednote{原文 voir Saussure s’est 如此；在此句法中应为 voir que Saussure s’est。}

\minititle{2. L’arbitraire selon Saussure}
Pour Saussure, le signe linguistique est et conventionnel et arbitraire. S’il existait un lien naturel entre le signifiant et le signifié, alors toutes les langues du monde utiliseraient le même signifiant. Il n’y aurait, comme le dit Leibnitz, « qu’une langue parmi les hommes » :
\begin{lecture}{}
Le lien unissant le signifiant au signifié est arbitraire, ou encore, puisque nous entendons par signe le total résultant de l’association d’un signifiant à un signifié, nous pouvons dire plus simplement le signe linguistique est arbitraire. Ainsi l’idée de « sœur » n’est liée par aucun rapport intérieur avec la suite de sons s - ö - r qui lui sert de signifiant ; il pourrait être aussi bien représenté par n’importe quelle autre : à preuve les différences entre les langues et l’existence même de langues différentes : le signifié « bœuf » a pour signifiant b - ö - f » d’un côté de la frontière, et o - k - s (Ochs) de l’autre.
\end{lecture}
L’arbitraire du signe chez Saussure est une autre façon de dire que signifiants n’ont aucun lien naturel avec la réalité à laquelle ils se rapportent.\ednote{此处把所指 signifié 换成了外部现实 réalité。CLG 的任意性原则直接涉及能指与所指的联系，不能不加区别地改说成符号与外部对象的关系。见\ecite{38}，第一部分第 I 章。} Du point du vue philosophique, Saussure se situe donc dans le camp du nominalisme. Mais s’il est vrai que le lien entre le signifiant et le signifié est arbitraire, que les « législateurs » des mots (c’est-à-dire la communauté linguistique) sont\origcont{184}{170} libres de nommer les objets comme bon leur semble, il n’en est pas moins vrai qu’une fois la nomination effectuée, le signifiant conventionnel s’impose au locuteur individuel. Les mots sont arbitraires, mais leur utilisation pour un groupe (langue) ou un individu (parole) ne l’est pas.
\begin{lecture}{}
Le mot arbitraire appelle aussi une remarque. Il ne doit pas donner l’idée que le signifiant dépend du libre choix du sujet parlant (on verra plus bas qu’il n’est pas au pouvoir de l’individu de rien changer à un signe une fois établi dans un groupe linguistique) ; nous voulons dire qu’il est immotivé, c’est-à-dire arbitraire par rapport au signifié, avec lequel il n’a aucune attache naturelle dans la réalité. [...] Le langage découpe simultanément un signifiant dans la masse informe des sons et un signifié dans la masse informe des concepts. Le rapport entre le signifiant et le signifié est arbitraire et immotivé : rien, a priori, ne justifie, en français par exemple, qu’à la suite de phonèmes {\BBSixGreek [a-ʁ-b-ʁ]} (le signifiant, en l’occurrence, du signe « arbre ») on associe le concept d’« arbre » (le signifié). Aucun raisonnement ne peut conduire à préférer {\BBSixGreek [bœf]} à {\BBSixGreek [ɔks]} pour signifier le concept de « bœuf ».
\end{lecture}
Dans l’extrait ci-dessus, le mot « informe » devrait attirer votre attention. Pour Saussure, la langue structure, organise, met de l’ordre dans la réalité informe qui nous entoure. D’où la célèbre formule de Saussure : « la langue est une forme, et non une substance ». On retrouve en filigrane cette idée dans l’extrait suivant, dans lequel l’adjectif « amorphe » a le même sens que l’adjectif « informe » :
\begin{lecture}{}
Psychologiquement, abstraction faite de son expression par les mots, notre pensée n’est qu’une masse amorphe et indistincte. Philosophes et linguistes se sont toujours accordés à reconnaître que, sans le secours des signes, nous serions incapables de distinguer deux idées d’une façon claire et constante. Prise en elle-même la pensée est comme une nébuleuse où rien n’est nécessairement délimité. Il n’y a pas d’idées préétablies, et rien n’est distinct avant l’apparition de la langue.
\end{lecture}
Autrement dit, Saussure s’inscrit dans la lignée de Hegel pour qui, nous l’avons vu, « c’est dans les mots que nous pensons ».

Pour Saussure, si le signe est arbitraire, tous les signes n’ont cependant pas le même degré d’arbitraire. Les symboles, on l’a vu avec Peirce, sont conventionnels. D’autres signes, en revanche, sont basés sur la ressemblance, et sont donc « moins conventionnels » :\ednote{Peirce 的 symbole 与 Saussure 在下引文字中的 symbole 不是同一个术语定义：前者以规律或习惯性联系界定，后者特别指出仍有一定自然联系。不能仅凭相同词形把两套定义接成同一分类；亦应与 CLG 的“相对任意性／相对理据性”区分。见\ecite{37}、\ecite{38}。}
\begin{lecture}{}
On s’est servi du mot symbole pour désigner le signe linguistique, ou plus exactement ce que nous appelons le signifiant, il y a des inconvénients à l’admettre, justement à cause de notre premier principe. Le symbole a pour caractère de n’être jamais tout à fait arbitraire ; il n’est pas vide, il y a un rudiment de lien naturel entre le signifiant et le signifié. Le symbole de la justice, la balance, ne pourrait pas être remplacé par n’importe quoi, un char, par exemple.
\end{lecture}

\origpage{185}{171}
\minititle{3. Le cas des onomatopées}
Saussure mentionne un autre type de signe dont la particularité est de ne pas être totalement arbitraire : il s’agit des onomatopées.
\begin{lecture}{}
On pourrait s’appuyer sur les onomatopées pour dire que le choix du signifiant n’est pas toujours arbitraire. Mais elles ne sont jamais des éléments organiques d’un système linguistique. Leur nombre est d’ailleurs bien moins grand qu’on ne le croit. Des mots comme fouet ou glas peuvent frapper certaines oreilles par une sonorité suggestive ; mais pour voir qu’ils n’ont pas ce caractère dès l’origine, il suffit de remonter à leurs formes latines (fouet dérivé de fāgus « hêtre », glas = classicum) ; la qualité de leurs sons actuels, ou plutôt celle qu’on leur attribue, est un résultat fortuit de l’évolution phonétique. Quant aux onomatopées authentiques (celles du type glou-glou, tic-tac, etc.), non seulement elles sont peu nombreuses, mais leur choix est déjà en quelque mesure arbitraire, puisqu’elles ne sont que l’imitation approximative et déjà à demi conventionnelle de certains bruits (comparez le français ouaoua et l’allemand wauwau). En outre, une fois introduites dans la langue, elles sont plus ou moins entraînées dans l’évolution phonétique, morphologique, etc. que subissent les autres mots (cf. pigeon, du latin vulgaire pīpiō, dérivé lui-même d’une onomatopée) : preuve évidente qu’elles ont perdu quelque chose de leur caractère premier pour revêtir celui du signe linguistique en général, qui est immotivé.
\end{lecture}
Un exemple permet d’illustrer cela. Prenons le cri du coq, tel que retranscrit dans différentes langues. On constate que le « cocorico » français est une imitation du son entendu. Cette imitation est en outre relativement proche dans la majorité des langues, comme on peut le constater ci-dessous :
\begin{center}\small
\begin{tabularx}{\linewidth}{@{}XX@{}}
$\diamond$ Angleterre : Cock-a-doodle-doo. & $\diamond$ France : Cocorico.\\[.5em]
$\diamond$ Belgique (Flamand) : Kukeleku. & $\diamond$ Grèce : Kiriki.\\[.5em]
$\diamond$ Belgique (Wallon) : Coucoulou Djoû. & $\diamond$ Hollande : Kukeleku.\\[.5em]
$\diamond$ Catalan : Kikiriki. & $\diamond$ Hongrie : Kukurikuuu.\\[.5em]
$\diamond$ Chine (Mandarin) : Gwo gwo. & $\diamond$ Inde : Kukru ku.\\[.5em]
$\diamond$ Corée : Kko-kki-yo. & $\diamond$ Indonésie : Kikeriku.\\[.5em]
$\diamond$ Croatie : Ku-ku-ri-ku. & $\diamond$ Italie : Chicchirichi.\\[.5em]
$\diamond$ Danemark : Kykkeliky. & $\diamond$ Luxembourg : kirikiki.\\[.5em]
$\diamond$ Espagne : Kirikiki. & $\diamond$ Norvège : Kykeliky.\\[.5em]
$\diamond$ Estonie : Kikerikii. & $\diamond$ Pays de Galles : coc-a-dwdl-dŵ.\\[.5em]
$\diamond$ Finlande : Kukkokiekuu. & $\diamond$ Pologne : Kukuryku.\\
\end{tabularx}
\end{center}
\origpage{186}{172}
\begin{center}\small
\begin{tabularx}{\linewidth}{@{}XX@{}}
$\diamond$ Portugal : Cocorococo. & $\diamond$ Thailande : Ake-e-ake-ake.\\[.5em]
$\diamond$ Roumanie : Cucurigu. & $\diamond$ Turquie : Kuk-kurri-kuuu.\\[.5em]
$\diamond$ Russie : Ku-ka-re-ku. & $\diamond$ Tchéquie : Kykyryký.\\[.5em]
$\diamond$ Suisse : Kuckeliku. & $\diamond$ Japon : Kokekokko\\
\end{tabularx}
\end{center}
Un autre exemple est le cri du canard, basé lui aussi sur l’imitation du son entendu :
\begin{itemize}
\item français : couin-couin
\item anglais : quack-quack
\item allemand : pack-pack
\item danois : rap-rap
\item hongrois : hap-hap
\end{itemize}
La réponse de Saussure est double : d’une part, on l’a vu dans l’extrait un peu plus haut, les onomatopées sont peu nombreuses : « En résumé, les onomatopées et les exclamations sont d’importance secondaire, et leur origine symbolique en partie contestable ». D’autre part, malgré des réelles ressemblances entre elles, elles varient d’un pays à l’autre, parfois de façon considérable. Ainsi, le « coin-coin » français est diamétralement différent du « hap-hap » hongrois. De même, le « Ake-e-ake-ake » thaïlandais est très différent du « chicchirichi » italien. Ceci révèle la part d’arbitraire plus ou moins grande des onomatopées, dont l’existence ne remet pas en cause le principe général d’arbitraire. C’est cela qui permet à Saussure de conclure que « le principe de l’arbitraire du signe n’est contesté par personne ».

\minititle{\textcircled{\scriptsize 3}\quad La linéarité du signe linguistique}
Après la dualité et l’arbitraire, le dernier caractère du signe linguistique (et peut-être le plus évident) est la linéarité.\ednote{CLG 的第二原则准确标题为“能指的线条性”（caractère linéaire du signifiant），首先限定于听觉能指；不宜把整个符号或一切语言模态都说成只能一维排列。下引首句保留了这一限定。见\ecite{38}，第一部分第 I 章 § 3。}
\begin{lecture}{}
Le signifiant, étant de nature auditive, se déroule dans le temps et a les caractères qu’il emprunte au temps : a) il représente une étendue et b) cette étendue est mesurable dans une seule dimension : c’est une ligne. Ce principe est évident, mais il semble qu’on ait toujours négligé de l’énoncer, sans doute parce qu’on l’a trouvé trop simple ; cependant il est fondamental et les conséquences en sont incalculables ; son importance est égale à celle de la première loi. Tout le mécanisme de la langue en dépend.
\end{lecture}
Ce principe est certes évident, mais il constitue une différence de taille avec les signes non linguistiques de la sémiologie, comme les signaux routiers, qui peuvent se déployer simultanément dans le temps et dans l’espace.

\origpage{187}{173}
\begin{lecture}{}
Par opposition aux signifiants visuels (signaux maritimes, etc.) qui peuvent offrir des complications simultanées sur plusieurs dimensions, les signifiants acoustiques ne disposent que de la ligne du temps ; leurs éléments se présentent l’un après l’autre ; ils forment une chaîne. Le signifiant se présente de façon linéaire dans l’axe du temps. Il nous faut du temps pour prononcer un mot, pour le réaliser de façon physique. De même, il y a un ordre qui est suivi lors de sa prononciation. Dans la réalisation du signifiant [wazo], il ne m’est pas permis de prononcer les sons dans un ordre différent de celui que nous avons ci-haut si je veux que les autres locuteurs me comprennent. Les signes forment donc une successivité et non une simultanéité. Par opposition, les signes routiers peuvent se substituer : obligation de tourner et tourner à gauche.
\end{lecture}
Les conséquences de cette linéarité sont importantes et seront détaillées un peu plus bas.

\minititle{\textcircled{\scriptsize 4}\quad La notion de valeur linguistique}
Pour Saussure, la langue est un système clos de signes. Dans ce système, tout signe est défini par rapport aux autres, par pure différence (négativement), et non par ses caractéristiques propres (positivement) :\ednote{这一说法需要 CLG 中的限定：能指或所指各自具有差异性，但“符号作为整体”被明确称为在自身秩序中的积极事实（fait positif）。不能删去这一层区别而一概否认整个符号的积极性质。见\ecite{38}，第二部分第 IV 章 § 4。}
\begin{lecture}{}
Placé dans un syntagme, un terme n’acquiert sa valeur que parce qu’il est opposé à ce qui précède ou ce qui suit, ou à tous les deux.
\end{lecture}

\minititle{1. Chez Patañjali}
Rappelons que Saussure était un spécialiste du sanskrit. Or, on retrouve très tôt chez les Hindous le débat sur la valeur linguistique. Ainsi, au 1\textsuperscript{er} siècle avant J.-C., \textbf{Patañjali} (le commentateur de Panini) défendait déjà l’idée que le signe se définit par ce qui l’oppose aux autres signes. À la question « qu’est-ce que le mot « \textit{gauḥ} » (bœuf) », il répondait que ce mot est ce qui fait naître, quand on le prononce, la notion de ce qui est muni d’une queue, de cornes, de pieds fourchus, etc. En d’autres termes, un concept se définit par une série de traits l’opposant à d’autres concepts. Encore un exemple de ce que la linguistique occidentale moderne et contemporaine doit à la grammaire de Panini.

\minititle{2. Chez Saussure}
Contrairement à une idée répandue, la langue n’est pas un répertoire de mots (Saussure parle de « nomenclature ») qui désigne les choses ou des concepts préexistants en leur collant des étiquettes. Si c’était le cas, les mots d’une langue auraient toujours leur correspondant exact dans une autre langue. Or, ce n’est pas le cas, et ceci va conduire Saussure à distinguer « signification » et « valeur ». Pour illustrer cela, il prend l’exemple de « mouton » en français et « \textit{sheep} » en anglais, qui ont la même signification mais pas la même valeur. En effet, l’anglais distingue \textit{sheep} (l’animal vivant) de sa viande (« \textit{mutton} »), ce qui n’est pas le cas en français. Ainsi, le signifié est un concept défini négativement du fait de l’existence ou de l’absence d’autres concepts qui lui sont
\origcont{188}{174} opposables. À l’intérieur d’une même langue, tous les mots qui expriment des idées voisines se limitent réciproquement. Des synonymes comme « redouter », « craindre », « avoir peur », n’ont de valeur que par leur opposition.

Signifié et signifiant se définissent négativement :
\begin{lecture}{}
Si la partie conceptuelle de la valeur est constituée uniquement par des rapports et des différences avec les autres termes de la langue, on peut en dire autant de sa partie matérielle. Ce qui importe dans le mot, ce n’est pas le son lui-même, mais les différences phoniques qui permettent de distinguer ce mot de tous les autres, car ce sont elles qui portent la signification.
\end{lecture}
La conclusion de Saussure est que le sens n’est pas à l’intérieur du mot mais dépend de l’environnement de ce dernier, c’est-à-dire « ce qu’il y a autour de lui dans les autres signes ». Par exemple la différence de sens entre « un homme grand » et « un grand homme » a été causée par le changement de l’environnement du mot « homme » :
\begin{lecture}{}
Tout ce qui précède revient à dire que dans la langue il n’y a que des différences. Ce qu’il y a d’idée ou de matière phonique dans un signe importe moins que ce qu’il y a autour de lui dans les autres signes. La preuve en est que la valeur d’un terme peut être modifiée sans qu’on touche ni à son sens ni à ses sons, mais seulement par le fait que tel autre terme voisin aura subi une modification. [...] Autrement dit, la langue est une forme et non une substance.
\end{lecture}

\minititle{\textcircled{\scriptsize 5}\quad Immutabilité du signe et mutabilité}
Pour Saussure, « il n’est pas au pouvoir de l’individu de rien changer à un signe une fois établi ».

On pourrait donc en conclure que le signe linguistique, une fois établi, est immuable. En fait, ce n’est pas si simple que cela : un peu comme la flèche qui, dans le poème de Paul Valéry, « vole et ne vole pas », comme « Achille immobile à grand pas », le signe linguistique change... et ne change pas !

\minititle{1. Immutabilité du signe}
Prenons un exemple très simple : Pour écrire ce livre, j’ai utilisé un ordinateur. Lorsque le mot « ordinateur » a été créé, on ne m’a pas demandé mon avis. D’autres l’ont décidé pour moi et avant moi. Et d’ailleurs, je n’ai pas le pouvoir d’utiliser un autre terme pour le désigner. C’est la même chose pour tous les mots que nous employons. Lorsque l’on apprend une langue, les mots nous sont imposés, par une sorte de loi (le « législateur » de Platon), de contrat : les mots de la langue échappent à notre volonté. Une langue, c’est quelque chose que l’on subit :

\origpage{189}{175}
\begin{lecture}{}
Si par rapport à l’idée qu’il représente, le signifiant apparaît comme librement choisi, en revanche, par rapport à la communauté linguistique qui l’emploie, il n’est pas libre, il est imposé. La masse sociale n’est point consultée, et le signifiant choisi par la langue, ne pourrait pas être remplacé par un autre. [...] La langue ne peut donc plus être assimilée à un contrat pur et simple, et c’est justement de ce côté que le signe linguistique est particulièrement intéressant à étudier ; car si l’on veut démontrer que la loi admise dans une collectivité est une chose que l’on subit, et non une règle librement consentie, c’est bien la langue qui en offre la preuve la plus éclatante. Voyons donc comment le signe linguistique échappe à notre volonté, et tirons ensuite les conséquences importantes qui découlent de ce phénomène.
\end{lecture}
La langue est un héritage, un produit achevé que l’on doit accepter comme tel, et c’est pour cela que le signe est immuable : « Nous disons homme et chien parce qu’avant nous on a dit homme et chien ». Quand ce contrat entre les mots et les idées a-t-il été passé ? Nous n’en savons rien, nous ne pouvons pas le savoir, donc cette question est inutile :
\begin{lecture}{}
À n’importe quelle époque, et si haut que nous remontions, la langue apparaît toujours comme un héritage de l’époque précédente. L’acte par lequel, à un moment donné, les noms seraient distribués aux choses, par lequel un contrat serait passé entre les concepts et les images acoustiques — cet acte, nous pouvons le concevoir, mais il n’a jamais été constaté. L’idée que les choses auraient pu se passer ainsi nous est suggérée par notre sentiment très vif de l’arbitraire du signe. En fait, aucune société ne connaît et n’a jamais connu la langue autrement que comme un produit hérité des générations précédentes et à prendre tel quel. C’est pourquoi la question de l’origine du langage n’a pas l’importance qu’on lui attribue généralement. Ce n’est pas même une question à poser ; le seul objet réel de la linguistique, c’est la vie normale et régulière d’un idiome déjà constitué. Un état de langue donné est toujours le produit de facteurs historiques, et ce sont ces facteurs qui expliquent pourquoi le signe est immuable, c’est-à-dire résiste à toute substitution arbitraire.
\end{lecture}
La langue est un produit social : elle est donc attachée au poids de la collectivité, et cela lui donne son caractère de fixité. Elle est aussi un héritage, et la « solidarité avec le passé » nous empêche de choisir. Et pourtant :
\begin{lecture}{}
Dire que la langue est un héritage n’explique rien si l’on ne va pas plus loin. Ne peut-on pas modifier d’un moment à l’autre des lois existantes et héritées ?
\end{lecture}
Saussure introduit un peu de liberté dans son propos :
\begin{lecture}{}
Il faut d’abord apprécier le plus ou moins de liberté dont jouissent les autres institutions ; on verra que pour chacune d’elles il y a une balance différente entre la tradition imposée et l’action libre de la société.
\end{lecture}

\origpage{190}{176}
\minititle{2. Mutabilité du signe}
C’est une particularité et un paradoxe du signe linguistique : les deux éléments qui confèrent à la langue sa fixité, à savoir les masses parlantes et le temps, sont les mêmes qui lui confèrent sa mutabilité.
\begin{lecture}{}
Le temps, qui assure la continuité de la langue, a un autre effet, en apparence contradictoire au premier : celui d’altérer plus ou moins rapidement les signes linguistiques et, en un certain sens, on peut parler à la fois de l’immutabilité et de la mutabilité du signe.
\end{lecture}
Il est important de noter ici l’emploi du verbe « altérer » : « l’infidélité au passé » n’est que relative, et cette altération s’inscrit dans une continuité. La question est de savoir quelle forme peut prendre cette altération :
\begin{lecture}{}
L’altération dans le temps prend diverses formes, dont chacune fournirait la matière d’un important chapitre de linguistique. Sans entrer dans le détail, voici ce qu’il est important de dégager. Tout d’abord, ne nous méprenons pas sur le sens attaché ici au mot altération. Il pourrait faire croire qu’il s’agit spécialement des changements phonétiques subis par le signifiant, ou bien des changements de sens qui atteignent les concepts signifiés. Cette vue serait insuffisante. Quels que soient les facteurs d’altération, qu’ils agissent isolément ou combinés, ils aboutissent toujours à un déplacement du rapport entre le signifié et le signifiant. [...] Une langue est radicalement impuissante à se défendre contre les facteurs qui déplacent d’instant en instant le rapport du signifié au signifiant.
\end{lecture}
On voit bien ici que la langue est soumise à deux influences contradictoires, qui font que l’on peut parler à la fois de mutabilité et d’immutabilité :
\begin{lecture}{}
Rien de plus complexe : située à la fois dans la masse sociale et dans le temps, personne ne peut rien y changer, et, d’autre part, l’arbitraire de ses signes entraîne théoriquement la liberté d’établir n’importe quel rapport entre la matière phonique et les idées. Il en résulte que ces deux éléments unis dans les signes gardent chacun leur vie propre dans une proportion inconnue ailleurs, et que la langue s’altère, ou plutôt évolue, sous l’influence de tous les agents qui peuvent atteindre soit les sons, soit les sens. Cette évolution est fatale ; il n’y a pas d’exemple d’une langue qui y résiste.
\end{lecture}

\subsection{Les dichotomies du CLG}
Comparée à la lecture de Peirce, la lecture du CLG de Saussure est une promenade de santé. La raison principale tient à la structure originale du livre bâti sur une série d’oppositions binaires de concepts, les dichotomies.

\origpage{191}{177}
\minititle{\textcircled{\scriptsize 1}\quad Linguistique interne et linguistique externe}
Pour Saussure, « la linguistique a pour unique et véritable objet la langue envisagée en elle-même et pour elle-même ». Sa linguistique est donc immanente. Comme le signe saussurien est dyadique, le référent est rejeté à l’extérieur du champ d’étude. Plus généralement, tout ce qu’il appelle « éléments externes à la langue », c’est-à-dire tout ce qui est extralinguistique, est rejeté hors du champ de la linguistique.\ednote{此处把研究对象的区分说成完全排除，过于绝对。CLG“导论”第 V 章明确承认外部语言学研究的重要性，但主张不必借助外部因素才能认识语言内部组织；不宜把这一方法论区分等同于外部语言学不属于语言研究。见\ecite{38}。} Il prône donc une linguistique interne, structurale, dont principales les branches sont :\ednote{语序误，应为 dont les principales branches sont。下列“语用学与言语行为”一项是本教材的现代编排，不宜读作 CLG 原书给出的分科表。}
\begin{itemize}
\item la phonologie : l’étude du rôle des sons dans le système linguistique
\item la morphologie : l’étude de la structure grammaticale des mots
\item la lexicologie : l’étude de vocabulaire composant le lexique d’une langue
\item la sémantique : l’étude de la signification
\item la syntaxe : l’étude des combinaisons et des relations des éléments de la phrase
\item la pragmatique et l’énonciation : l’étude de la production et de la reconnaissance langagière par des énonciateurs dans un contexte donné
\end{itemize}
Dès le chapitre suivant, dans la deuxième partie de notre livre, nous allons étudier chacune de ces branches. Quant aux branches de la linguistique externe, elles concernent le locuteur ou les masses parlantes, et comprennent par exemple :
\begin{itemize}
\item la sociolinguistique, qui étudie des comportements langagiers collectifs.
\item la géolinguistique, étude linguistique des variantes régionales d’une langue.
\item la neurolinguistique, qui étudie la corrélation entre la structure langagière et la structure neurologique du locuteur.
\end{itemize}
Ces branches seront quant à elles présentées dans la troisième partie de notre livre, consacrée aux développements récents des sciences du langage.

\minititle{\textcircled{\scriptsize 2}\quad Linguistique de la langue et linguistique de la parole}
Dans l’extrait suivant, Saussure établit la différence entre langue et langage, nous y retrouvons des idées bien familières :
\begin{lecture}{}
Mais qu’est-ce que la langue ? Pour nous elle ne se confond pas avec le langage ; elle n’en est qu’une partie déterminée, essentielle, il est vrai. C’est à la fois un produit social de la faculté du langage et un ensemble de conventions nécessaires, adoptées par le corps social pour permettre l’exercice de cette faculté chez les individus.
\end{lecture}

\origpage{192}{178}
\minititle{1. Le langage, hétéroclite et insaisissable}
Le langage, en raison de sa nature hybride, insaisissable, et inconnaissable, ne se prête à aucune classification. C’est un problème car lorsque l’on veut établir une science, son objet doit être uniforme, homogène et clairement défini :
\begin{lecture}{}
Pris dans son tout, le langage est multiforme et hétéroclite à cheval sur plusieurs domaines, à la fois physique, physiologique et psychique, il appartient encore au domaine individuel et au domaine social ; il ne se laisse classer dans aucune catégorie des faits humains, parce qu’on ne sait comment dégager son unité. La langue, au contraire, est un tout en soi et un principe de classification. Dès que nous lui donnons la première place parmi les faits de langage, nous introduisons un ordre naturel dans un ensemble qui ne se prête à aucune autre classification. [...]
\end{lecture}

\minititle{2. La langue, principe de classification}
Saussure va donc choisir la langue, « objet bien défini », comme objet de la linguistique et la placer au centre des faits de langage.
\begin{lecture}{}
La langue est un objet bien défini dans l’ensemble hétéroclite des faits de langage. [...] Elle est la partie sociale du langage, extérieure à l’individu qui à lui seul ne peut ni la créer ni la modifier ; elle n’existe qu’en vertu d’une sorte de contrat passé entre les membres de la communauté.
\end{lecture}
Le langage est la somme de la langue et de la parole. Après avoir distingué le langage et la langue, Saussure va faire une autre distinction déterminante.

\minititle{3. La parole, partie individuelle du langage}
\begin{lecture}{}
L’étude du langage comporte donc deux parties l’une, essentielle, a pour objet la langue, qui est sociale dans son essence et indépendante de l’individu ; cette étude est uniquement psychique ; l’autre, secondaire, a pour objet la partie individuelle du langage, c’est-à-dire la parole y compris la phonation, elle est psycho-physique.
\end{lecture}
Bien sûr, de même qu’il est impossible de séparer signifiant et signifié, il est impossible de séparer langue et parole :
\begin{lecture}{}
Ces deux objets sont étroitement liés et se supposent l’un l’autre ; la langue est nécessaire pour que la parole soit intelligible et produise tous ses effets ; mais celle-ci est nécessaire pour que la langue s’établisse ; historiquement, le fait de parole précède toujours.
\end{lecture}
Sans parole, pas de langue, et sans langue pas de parole. Il y a donc, dit Saussure, « interdépendance de la langue et de la parole ; celle-là est à la fois l’instrument et le produit de celle-ci ».

\origpage{193}{179}
\minititle{4. Langue ou parole, il faut choisir}
Leur interdépendance n’occulte pas leurs différences : dès lors, le linguiste doit choisir deux voies distinctes. Pour Saussure, le choix est clair : c’est la langue.
\begin{lecture}{}
Pour toutes ces raisons, il serait chimérique de réunir sous un même point de vue la langue et la parole. Le tout global du langage est inconnaissable, parce qu’il n’est pas homogène, tandis que la distinction et la subordination proposées éclairent tout. Telle est la première bifurcation qu’on rencontre dès qu’on cherche à faire la théorie du langage. Il faut choisir entre deux routes qu’il est impossible de prendre en même temps ; elles doivent être suivies séparément. On peut à la rigueur conserver le nom de linguistique à chacune de ces deux disciplines et parler d’une linguistique de la parole. Mais il ne faudra pas la confondre avec la linguistique proprement dite, celle dont la langue est l’unique objet.
\end{lecture}

\minititle{\textcircled{\scriptsize 3}\quad Linguistique synchronique et linguistique diachronique}
Au-delà du choix entre linguistique de la langue et linguistique de la parole, le linguiste doit faire un autre choix : étudier l’évolution des signes au cours du temps, ou bien les rapports entre les signes à un instant donné. Dans le premier cas, sa démarche est diachronique, dans le second elle est synchronique. C’est dans l’étude de ce second aspect que Saussure a particulièrement innové. Selon lui, la perspective diachronique doit être étudiée, certes, mais elle ne permet pas de rendre compte du fait que la langue est un système. Elle prend en effet uniquement en compte les modifications au cours du temps : l’approche synchronique montre que la signification des signes dépend de la structure de l’ensemble de la langue.

\minititle{1. La synchronie}
Pour définir ces deux notions, voici quelques extraits des \textit{Leçons de linguistique} que le linguiste français \textbf{Gustave Guillaume} (1883-1960) a données entre 1945-1946.
\begin{lecture}{}
Ferdinand de Saussure fait la distinction d’une linguistique évolutive qui observe les faits de langue et de discours sur l’axe des successivités, et d’une linguistique statique dont la vocation est d’observer les mêmes faits mais d’un point de vue différent, sur l’axe des états. Quant à l’observation des faits de langue et des faits de discours sur l’axe des états, Saussure considère qu’elle procède d’un point de vue qu’il nomme la synchronie.

La linguistique descriptive, et tout particulièrement celle qui prend pour objet la systématisation intérieure des langues, est une discipline relevant de la synchronie. Il lui faut en effet observer les êtres de langue, quelle qu’en soit l’espèce, sous l’angle de leur capacité de fonctionnement à un moment donné dont on ne sort pas.
\end{lecture}

\origpage{194}{180}
\begin{lecture}{}
La linguistique synchronique suppose, pour tout instant historique considéré, la vue du système que constitue la langue, et la connaissance que cette linguistique recherche est à celle de la position, dans le système aperçu, de chacune des parties constitutives.\ednote{引文“est à celle”在本句中语法不通，疑为 est celle；未取得此引文所据 Gustave Guillaume 课程版次及页码，故只列拟校，不改正文。} L’examen des parties, fait attentivement, conduit à une perception de la structure d’ensemble, et la connaissance de celle-ci conduit à son tour à une vision exacte des parties composantes et de leur relativité réciproque dans le tout systématique qu’elles recomposent par leur réunion.
\end{lecture}

\minititle{2. La diachronie}
Voici à présent comment Gustave Guillaume définit la diachronie :
\begin{lecture}{}
L’observation des faits de discours et des faits de langue sur l’axe des successivités procède d’un point de vue que Saussure désigne d’un mot heureusement choisi : la diachronie. La grammaire historique et, à côté d’elle, usant de moyens différents, la grammaire comparative ont des disciplines relevant de la diachronie.\ednote{原文“ont des disciplines”疑漏 s，应为 sont des disciplines。仍需与所据课程版本核对。}
\end{lecture}
Bien que formé au comparatisme et qu’il ait été professeur de grammaire comparative, Saussure a opté pour la synchronie. Pour lui, l’objet du linguiste, c’est un état de langue, un système à un moment donné.

\minititle{\textcircled{\scriptsize 4}\quad Rapports syntagmatiques et rapports paradigmatiques}
\begin{lecture}{}
Pour Saussure, dans la langue, tout repose sur des rapports, et les mots entretiennent entre eux deux types de rapports : les rapports syntagmatiques et les rapports paradigmatiques. Ces deux rapports correspondent à deux formes différentes mais complémentaires de l’activité mentale de l’homme, et « sont toutes deux indispensables à la vie de la langue ».
\end{lecture}
\ednote{本框不是 CLG 的逐字引文。CLG 第二部分第 V 章使用 rapports associatifs，与 syntagmatiques 对举；“paradigmatiques”是这里采用的后起术语。关联关系也不只限于同一语法位置可替换的项目，后文举 enseigner、armement 等即显示其范围更广。见\ecite{38}。}

\minititle{1. Les rapports syntagmatiques}
Cette notion, expliquée ci-dessous par Saussure, paraît simple en apparence mais ses conséquences sont importantes :
\begin{lecture}{}
D’une part, dans le discours, les mots contractent entre eux, en vertu de leur enchaînement, des rapports fondés sur le caractère linéaire de la langue, qui exclut la possibilité de prononcer deux éléments à la fois. Ceux-ci se rangent les uns à la suite des autres sur la chaîne de la parole. Ces combinaisons qui ont pour support l’étendue peuvent être appelées syntagmes. Le syntagme se compose donc toujours de deux ou plusieurs unités consécutives.
\end{lecture}

\minititle{1) Illustration \# 1}
Plus de clarté, illustrons ce concept par un exemple :\ednote{Plus de clarté 在此缺少引介介词，宜为 Pour plus de clarté。}

\noindent\textbf{Ex :} Le vieux professeur écrit un livre.

\origpage{195}{181}
Si on lit cette phrase à haute voix, on produit une séquence de sons se succédant chronologiquement, selon un ordre imposé. L’ensemble des règles qui déterminent l’ordre selon lequel certains éléments peuvent être mis en séquence est la syntaxe. Un groupe d’éléments doté d’un sens propre est un syntagme. Quant à l’enchaînement des éléments de l’énoncé selon les règles de la syntaxe, on l’appelle concaténation (du latin \textit{catena}, la chaîne). On peut schématiser la phrase comme suit :

\noindent Le $\to$ vieux $\to$ professeur $\to$ écrit $\to$ un $\to$ livre

Ici, le signe $\to$ symbolise les relations syntagmatiques : les éléments ne succèdent dans la phrase selon certaines contraintes de syntaxe, sur un axe des successivités.\ednote{原文 ne succèdent 应为 se succèdent。} On peut traduire cette relation syntagmatique par un « et » : on peut trouver un élément A (déterminant) \textit{et} un élément B (adjectif) \textit{et} un élément C (nom), etc. La séquence suivante, en revanche, n’est pas possible en français :

\noindent *vieux $\to$ le $\to$ livre $\to$ le $\to$ écrit $\to$ professeur

Cette forme est agrammaticale car elle enfreint les lois de la syntaxe : le français est, rappelons-le, une langue SVO.

Nous pouvons visualiser les rapports syntagmatiques dans le tableau ci-dessous : on y voit la division de la phrase en deux syntagmes principaux : le syntagme nominal correspond au sujet, et le syntagme prédicatif correspond au prédicat. Le syntagme prédicatif se divise à son tour en syntagme verbal et en syntagme nominal. Chaque syntagme se décompose enfin en éléments simples placés sur l’axe syntagmatique.

\begin{center}\small
\renewcommand{\arraystretch}{1.35}
\begin{tabularx}{\linewidth}{|*{6}{>{\centering\arraybackslash}X|}}
\hline
\multicolumn{6}{|c|}{\textbf{Phrase}}\\\hline
\multicolumn{3}{|c|}{Syntagme Nominal}&\multicolumn{3}{c|}{Syntagme Prédicatif}\\\hline
\multicolumn{3}{|c|}{}&Syntagme Verbal&\multicolumn{2}{c|}{Syntagme Nominal}\\\hline
déterminant&adjectif&nom&verbe&déterminant&nom\\\hline
Le&vieux&professeur&écrit&un&livre\\\hline
\end{tabularx}
\end{center}

\minititle{2) Illustration \# 2}
Pour illustrer les rapports syntagmatiques et l’importance de l’ordre des mots en français, je vous propose un peu de détente avec un extrait du \textit{Bourgeois Gentilhomme} (1670) de Molière. Monsieur Jourdain, riche mais illettré, sollicite l’aide de son maître de philosophie pour écrire un billet galant (une lettre d’amour) à la belle et jeune marquise dont il est épris :

\origpage{196}{182}
\begin{lecture}{Lecture : Le Bourgeois Gentilhomme (extrait)}
Monsieur Jourdain : Je voudrais donc lui mettre dans un billet : Belle Marquise, vos beaux yeux me font mourir d’amour ; mais je voudrais que ce fût mis d’une manière galante, que cela fût tourné gentiment.

Maître de philosophie : Mettre que les feux de ses yeux réduisent votre cœur en cendres ; que vous souffrez nuit et jour pour elle les violences d’un...

Monsieur Jourdain : Non, non, non, je ne veux point de tout cela ; je ne veux que ce que je vous ai dit : Belle Marquise, vos beaux yeux me font mourir d’amour.

Maître de philosophie : Il faut bien étendre un peu la chose.

Monsieur Jourdain : Non, vous dis-je, je ne veux que ces seules paroles là dans le billet ; mais tournées à la mode, bien arrangées comme il faut. Je vous prie de me dire un peu, pour voir, les diverses manières dont on peut les mettre.

Maître de philosophie : On les peut mettre premièrement comme vous avez dit : Belle Marquise, vos beaux yeux me font mourir d’amour. Ou bien : D’amour mourir me font, belle Marquise, vos beaux yeux. Ou bien : Vos yeux beaux d’amour me font, belle Marquise, mourir. Ou bien : Mourir vos beaux yeux, belle Marquise, d’amour me font.

Ou bien : Me font vos yeux beaux mourir, belle Marquise, d’amour.

Monsieur Jourdain : Mais de toutes ces façons-là, laquelle est la meilleure ?

Maître de philosophie : Celle que vous avez dite : Belle Marquise, vos beaux yeux me font mourir d’amour.

Monsieur Jourdain : Cependant je n’ai point étudié, et j’ai fait cela tout du premier coup. Je vous remercie de tout mon cœur [...]
\end{lecture}

\minititle{2. Les rapports paradigmatiques}
Pour les éléments qui appartiennent à un même paradigme, c’est-à-dire qui peuvent remplir la même fonction dans la phrase, Saussure postule un deuxième type de relation, appelée relation paradigmatique. Comme ils impliquent la possibilité d’une substitution, on peut traduire ces rapports syntagmatiques par un « ou »,\ednote{此处 syntagmatiques 与本段标题及“替换／或”的解释相矛盾，应为 paradigmatiques。} par opposition au « et » vu un peu plus haut dans les rapports syntagmatiques. À un endroit donné dans une phrase, on peut en effet trouver un élément A \textit{ou} un élément B \textit{ou} un élément C, etc. :

\begin{center}\small
\renewcommand{\arraystretch}{1.35}
\begin{tabularx}{\linewidth}{|*{6}{>{\centering\arraybackslash}X|}}
\hline
\multicolumn{6}{|c|}{\textbf{Phrase}}\\\hline
\multicolumn{3}{|c|}{Syntagme Nominal}&\multicolumn{3}{c|}{Syntagme Prédicatif}\\\hline
\multicolumn{3}{|c|}{}&Syntagme Verbal&\multicolumn{2}{c|}{Syntagme Nominal}\\\hline
déterminant&adjectif&nom&verbe&déterminant&nom\\\hline
La&jeune&avocate&regarde&le&dossier\\\hline
Mon&petit&frère&allume&la&télé\\\hline
Les&jeunes&Français&adorent&le&vin\\\hline
\end{tabularx}
\end{center}

\origpage{197}{183}
Cette relation paradigmatique est aussi appelée « associative » par Saussure, et chaque catégorie (déterminant, adjectif, nom, verbe, etc.) est un paradigme. L’axe paradigmatique est vertical et représente l’axe des potentialités, c’est-à-dire ce que le locuteur peut dire, des mots qu’il peut choisir. Parler, pour un individu, c’est donc d’abord faire un choix parmi tous les mots dont il dispose (« le trésor intérieur », pour Saussure), puis les disposer sur la chaîne parlée en respectant les contraintes de la syntaxe. Les éléments ainsi combinés sont par ailleurs associés chez le locuteur à d’autres éléments qui ont entre eux « quelque chose de commun » : « enseignement », par exemple est associé à « enseignant » par parenté, à « armement » ou « chargement » par suffixation, ou encore à « apprentissage » et « éducation » par analogie sémantique :
\begin{lecture}{}
Les mots offrant quelque chose de commun s’associent dans la mémoire, et il se forme ainsi des groupes au sein desquels règnent des rapports très divers. Ainsi le mot enseignement fera surgir inconsciemment devant l’esprit une foule d’autres mots (enseigner, renseigner, etc., ou bien armement, changement, etc., ou bien éducation, apprentissage) ; par un côté ou un autre, tous ont quelque chose de commun entre eux.

On voit que ces coordinations sont d’une tout autre espèce que les premières. Elles n’ont pas pour support l’étendue leur siège est dans le cerveau ; elles font partie de ce trésor intérieur qui constitue la langue chez chaque individu. Nous les appellerons rapports associatifs.
\end{lecture}
Saussure a ainsi défini les deux types fondamentaux de relations entre les éléments d’un système, avec deux différences toutefois : alors que les rapports syntagmatiques sont directement observables et sont la conséquence directe de la linéarité du signe linguistique, les rapports associatifs sont virtuels, sous-jacents.

Pour conclure sur ce point et sur ce chapitre, Saussure utilise une comparaison issue de l’architecture pour « rendre visibles » les rapports syntagmatiques et paradigmatiques (associatifs) :
\begin{lecture}{}
À ce double point de vue, une unité linguistique est comparable à une partie déterminée d’un édifice, une colonne par exemple ; celle-ci se trouve, d’une part, dans un certain rapport avec l’architrave qu’elle supporte ; cet agencement de deux unités également présentes dans l’espace fait penser au rapport syntagmatique ; d’autre part, si cette colonne est d’ordre dorique, elle évoque la comparaison mentale avec les autres ordres (ionique, corinthien, etc.), qui sont des éléments non présents dans l’espace le rapport est associatif.
\end{lecture}
C’est sur cette image que s’achève le sixième chapitre mais aussi la première partie de notre livre, qui nous a mené des grammairiens antiques à la naissance de la linguistique moderne. Chacun des six chapitres de la deuxième partie sera consacré à l’un des six domaines de la linguistique interne : étude des sons (phonétique articulatoire), des phonèmes (phonologie), des morphèmes (morphologie), des sèmes (sémantique), des phrases (syntaxe) et des actes de langages (pragmatique).\ednote{术语通常写作 actes de langage，此处 langages 的复数不当。前面“语义学／句法学”的列举次序与全书目录第 10 章 Syntaxe、第 11 章 Sémantique 相反；正文仍按底本保留。}

\origpage{198}{184}
\section*{Exercices}
\phantomsection\addcontentsline{toc}{section}{Exercices}
\markright{Chapitre 6 — Exercices}
\noindent\textbf{I. Dans cet extrait du \textit{Cours de Linguistique Générale}, de quelle manière la métaphore du jeu d’échec permet-elle de donner corps aux concepts de synchronie, de diachronie, d’arbitraire, d’immutabilité et de valeur linguistique ?}
\begin{lecture}{}
Une partie d’échecs est comme une réalisation artificielle de ce que la langue nous présente sous une forme naturelle. Voyons la chose de plus près.

D’abord un état du jeu correspond bien à un état de la langue. La valeur respective des pièces dépend de leur position sur l’échiquier, de même que dans la langue chaque terme a sa valeur par son opposition avec tous les autres termes.

En second lieu, le système n’est jamais que momentané ; il varie d’une position à l’autre. Il est vrai que les valeurs dépendent aussi et surtout d’une convention immuable, la règle du jeu, qui existe avant le début de la partie et persiste après chaque coup. Cette règle admise une fois pour toutes existe aussi en matière de langue ; ce sont les principes constants de la sémiologie. Enfin, pour passer d’un équilibre à l’autre, ou - selon notre terminologie - d’une synchronie à l’autre, le déplacement d’une pièce suffit ; il n’y a pas de remue-ménage général. Nous avons là le pendant du fait diachronique avec toutes ses particularités. En effet, dans une partie d’échecs, n’importe quelle position donnée a pour caractère singulier d’être affranchie de ses antécédents ; il est totalement indifférent qu’on y soit arrivé par une voie ou par une autre ; celui qui a suivi toute la partie n’a pas le plus léger avantage sur le curieux qui vient inspecter l’état du jeu au moment critique; pour décrire cette position, il est parfaitement inutile de rappeler ce qui vient de se passer dix secondes auparavant. Tout ceci s’applique également à la langue et consacre la distinction radicale du diachronique et du synchronique. La parole n’opère jamais que sur un état de langue, et les changements qui interviennent entre les états n’y ont eux-mêmes aucune place.
\end{lecture}

\noindent\textbf{II. Expression personnelle :}

Pensez-vous, comme Saussure, que la question de l’origine du langage n’est pas une question à poser ?

\noindent\textbf{III. Retrouvez le sens des onomatopées suivantes.}
\begin{center}\renewcommand{\arraystretch}{1.2}
\begin{tabularx}{.78\linewidth}{|*{2}{>{\centering\arraybackslash}X|}}
\hline Areu areu&Paf\\\hline
Atchoum&Pan\\\hline
Badaboum&Patati, Patata\\\hline
\end{tabularx}\end{center}

\origpage{199}{185}
\begin{center}\small\renewcommand{\arraystretch}{1.18}
\begin{longtable}{|>{\centering\arraybackslash}p{.36\linewidth}|>{\centering\arraybackslash}p{.36\linewidth}|}
\hline
Beurk&Patatras\\\hline
Bof&Pfff\\\hline
Bouh&Pin Pon\\\hline
Coin coin&Piou Piou\\\hline
Cocorico&Plop\\\hline
Coa&Plouf\\\hline
Cot cot&Pouah\\\hline
Cui-cui&Pouf\\\hline
Croâ&Pschitt\\\hline
Ding Dong, Dong&Ron ron\\\hline
Drelin Drelin&Raaah\\\hline
Dring&Ratatatata\\\hline
Gla gla&Rrrr\\\hline
Glou Glou&Snif snif\\\hline
Groin groin&Splash\\\hline
Guili-guili&Tagada\\\hline
Hic&Tic Tac\\\hline
Hi han&Tchou tchou\\\hline
Hop&Tuuut\\\hline
Miam Miam&Ta Tac Ta Toum\\\hline
Meuh&Vlan\\\hline
Ouin&Vroum Vroum\\\hline
Ouille&\\\hline
\end{longtable}\end{center}

\noindent\textbf{IV. Xun zi（荀子）a dit : « 名无固宜，约之以命，约定俗成谓之宜，异于约则谓之不宜。名无固实，约之以命实，约定俗成，谓之实名。 » De quels concepts et penseurs occidentaux étudiés dans ce chapitre peut-on la rapprocher ? Comment la traduiriez-vous en français ?}

\origpage{200}{186}
\section*{Pour aller plus loin}
\phantomsection\addcontentsline{toc}{section}{Pour aller plus loin}
\markright{Chapitre 6 — Pour aller plus loin}
\begin{small}
\setlength{\parindent}{0pt}\setlength{\parskip}{.55em}
BARTHES (Roland). \textit{Le Degré zéro de l’écriture}. Paris, Seuil, 1953.

BARTHES (Roland). \textit{Mythologies}. Paris, Seuil, 2010.

BENVENISTE (Émile). \textit{Problèmes de linguistique générale}. Paris, Gallimard, vol 1, 1966, vol 2, 1974.

DELEDALLE (Gérard). \textit{Théorie et pratique du signe, Introduction à la sémiotique de Charles S. Peirce}. Paris, Payot, 1979.

DOROTHÉE (Stéphane). À l’origine du signe, \textit{le latin signum}. Paris, L’Harmattan, 2006.

GUIRAUD (Pierre). \textit{La sémiologie}, coll. « Que sais-je ? ». Paris, P.U.F., 1971.

MORRIS (Charles W.). \textit{Foundations of the Theory of Signs (1938)}. Tr. fr. partielle : « Fondements de la théorie des signes », dans \textit{Langages}. Paris, Larousse, 1974, vol. 35.

MORRIS (Charles W.). \textit{Signs, Language and Behavior}. New York, Prentice Hall, 1946.

PEIRCE (Charles Sanders). Écrits sur le signe. Paris, Seuil, 1978.

PRIETO (Luis Jorge). \textit{Messages et signaux}. Paris, Presses universitaires de France, 1966.

SAUSSURE (Ferdinand de). \textit{Cours de linguistique générale}. Paris, Payot, 1916.
\end{small}

\begingroup\setstretch{1.05}\setlength{\parskip}{.18em}
\section*{Glossaire}
\phantomsection\addcontentsline{toc}{section}{Glossaire}
\markright{Chapitre 6 — Glossaire}
\gloss{diachronie 历时}{研究语言在一段时间内如何变化的研究方式。}
\gloss{dyade 二元，（对话的）二人组合}{互相交流的两个人。二人组合是交际网络中最小的组成部分。\ednote{此释义属于交际参与者的二人组合；本章 dyade 指能指与所指构成的二元关系，不是两个交谈者。应按本章语境释为“二元组／二元结构”。}}
\gloss{sémiotique/sémiologie 符号学，符号学的}{关于符号的理论，或指对用符号或信号进行交际的系统的分析。}

\origpage{201}{187}
\gloss{signal 信号}{信息的发送者将所要传达的信息转化为信号，由通讯渠道传到接收者，再还原为信息传到终点。}
\gloss{signe 符号}{语言学中，指代表事物的单词或其他短语。\ednote{这一定义比本章的用法窄：本章还讨论非语言符号；Saussure 的语言符号也并非仅指“代表事物的单词”，而是能指与所指的结合，且不能化约为名称与事物的配对。见\ecite{38}，第一部分第 I 章。}}
\gloss{signifiant 能指}{语言学符号的具体含义（如某一读音指代的形象或某一图片展示的符号），是所指的任意表现。\ednote{“能指”误释为“具体含义”，与下条“所指”混淆。在 Saussure 的术语中，能指是音响形象（image acoustique），所指是概念；音响形象也不等于外界的物理声波。本章所述任意性针对两者间的联系。见\ecite{38}，第一部分第 I 章。}}
\gloss{signification 表指关系，语法意义}{指符号所代表的事物，一个语言项目所具有的作为语言系统中用法的意义。\ednote{“符号所代表的事物”与意义概念混淆，也与本章关于 signifié／référent 的区分不符。signification 可依语境指意义、意指或表意关系，并不专指“语法意义”；不应把它等同于指称对象。参见本章原书 157、160、173 页及\ecite{38}。}}
\gloss{signifié 所指}{语言学符号的语意或表示的概念，由能指具体表现出来。}
\gloss{symbole 象征性符号}{通常指代用来表现某种特别意义的事物。}

\endgroup
